\section{VII. 형식주의(Formalism)와
규칙회의주의(Rule-Scepticism)}\label{vii.-uxd615uxc2dduxc8fcuxc758formalismuxc640-uxaddcuxce59uxd68cuxc758uxc8fcuxc758rule-scepticism}

\subsection{1. 법의 개방적 직조(The Open Texture of
Law)}\label{uxbc95uxc758-uxac1cuxbc29uxc801-uxc9c1uxc870the-open-texture-of-law}

어떤 대규모 집단에서든, 일반 규칙들, 기준들, 원칙들은 사회 통제의 주요
수단이어야 하며, 각 개인에게 별도로 내려지는 개별 지시들이 주요
수단이어서는 안 된다. 수많은 개인들이 추가 지시 없이도, 어떤 경우가
발생했을 때 자신들에게 일정한 행위가 요구된다는 것으로 이해할 수 있는
일반적 행위 기준들을 전달할 수 없다면, 우리가 현재 법으로 인정하는 것은
아무것도 존재할 수 없을 것이다. 따라서 법은 주로, 그러나 결코 전적으로는
아니게, 사람의 \emphksans{부류들(classes)}, 그리고 행위, 사물, 상황의
\emphksans{부류들(classes)}을 가리켜야 하며, 사회생활의 광범위한 영역에서
법이 성공적으로 작동하는 것은, 구체적인 행위들, 사물들, 상황들을 법이
만드는 일반적 분류들의 사례들로 식별할 수 있는 능력(capacity)이 널리
퍼져 있는 데 의존한다.

그러한 일반적 행위 기준들을, 그것들이 적용될 이어지는 경우들에 앞서
전달하기 위해 사용되어 온 두 가지 주요 장치가 있다. 처음 보기에는 서로
매우 달라 보이는 이 두 장치 가운데 하나는 일반적 분류어(classifying
words)를 최대한 사용하고, 다른 하나는 최소한 사용한다. 첫 번째 방식은
우리가 입법(legislation)이라고 부르는 것으로 대표되고, 두 번째 방식은
선례(precedent)로 대표된다. 우리는 다음의 단순한 비법적 사례들에서 이
둘의 구별되는 특징들을 볼 수 있다. 한 아버지는 교회에 가기 전에 아들에게
이렇게 말한다. “모든 남성과 소년은 교회에 들어갈 때 모자를 벗어야
한다.” 다른 아버지는 교회에 들어가면서 자기 모자를 벗으며 말한다.
“보아라. 이런 경우에는 이렇게 행동하는 것이 옳은 방식이다.”

예시에 의해 행위 기준들(standards of conduct)을 전달하거나 가르치는
것은, 우리의 단순한 사례보다 훨씬 더 정교한 여러 형태를 취할 수 있다.
만약 특정한 경우에 아이에게 교회에 들어갈 때 아버지가 한 행동을 옳은
행위의 본보기로 간주하라고 말하는 대신, 아버지가 \textbf{(p.~125)}
아이가 자신을 적절한 행위에 관한 권위로 간주하고 어떻게 행동해야
하는지를 배우기 위해 자신을 관찰할 것이라고 가정한다면(assumed), 우리의
사례는 법에서의 선례(precedent) 사용에 더 가까워질 것이다. 법에서의 선례
사용에 더 가까이 접근하려면, 아버지가 자기 자신과 타인들에 의해 전통적인
행위 기준들을 지지하는 사람으로 이해되고 있으며 새로운 기준들을 도입하는
것이 아니라고 상정해야 한다.

예시에 의한 전달은 그 모든 형태에서, ‘내가 하는 대로 하라’와 같은
얼마간의 일반적 언어 지시가 동반되더라도, 의도된 것이 무엇인지에 관하여,
전달하려는 사람이 스스로 명확히 염두에 둔 사항들에 대해서조차,
가능성들의 범위와 따라서 의문의 범위를 열어둘 수 있다. 수행의 어느
정도까지가 모방되어야 하는가? 모자를 벗을 때 오른손 대신 왼손을 쓰는
것이 문제가 되는가? 느리게 하는가 재빠르게 하는가가 문제가 되는가?
모자를 좌석 아래에 두는 것이 문제가 되는가? 교회 안에서 모자를 다시 쓰지
않는 것이 문제가 되는가? 이 모든 것은 아이가 스스로에게 물을 수 있는
보다 일반적인 질문들의 변형들이다. ‘내 행위가 그의 행위와 어떤 방식으로
유사해야 옳은가?’ ‘그의 행위에서 정확히 무엇이 나의 지침이 되어야
하는가?’ 예시를 이해하면서 아이는 그 예시의 어떤 측면들에 주목하고 다른
측면들에는 주목하지 않는다. 그렇게 할 때 아이는 상식과, 어른들이
중요하다고 생각하는 사물들과 목적들의 일반적 종류에 관한 지식, 그리고 그
경우(교회에 가는 것)의 일반적 성격과 그에 적절한 행위 종류에 대한 이해에
의해 인도된다.

예시들의 불확정성들(indeterminacies)과 대조적으로, 명시적인 일반 언어
형식(예: “모든 남성은 교회에 들어갈 때 모자를 벗어야 한다”)에 의한
일반 기준들의 전달은 분명하고, 신뢰할 수 있으며, 확실한 것처럼 보인다.
행위의 일반적 지침들로 받아들여져야 할 특징들은 여기서 말로 식별된다.
그것들은 구체적인 예시 속에 다른 것들과 함께 묻혀 남겨져 있지 않고,
언어적으로 끄집어내어져 있다. 아이는 다른 경우들에서 무엇을 해야 할지를
알기 위해 더 이상 무엇이 의도되었는지, 또는 무엇이 승인될 것인지를
추측할 필요가 없다. 아이는 자기 행위가 옳기 위해서는 그 예시와 어떤
방식으로 유사해야 하는지를 추측해야 하는 처지에 남겨져 있지 않다. 대신
아이에게는 앞으로 무엇을 해야 하는지, 그리고 언제 그것을 해야 하는지를
골라낼 수 있는 언어적 기술(description)이 있다. 그는 명확한 언어
용어들의 사례들을 식별하고, 구체적 사실들을 일반적 분류 표제들 아래에
‘포섭(subsume)’하며, 단순한 삼단논법적 결론을 도출하기만 하면 된다.
\textbf{(p.~126)} 그는 자기 위험 부담 아래 선택하거나 추가적인 권위 있는
지침을 구해야 하는 선택지 앞에 놓여 있지 않다. 그는 자기 혼자서 자기에게
적용할 수 있는 규칙(rule)을 가지고 있다.

20세기의 법리학의 많은 부분은, 권위 있는 예시(선례(precedent))를 통한
전달의 불확실성과 권위 있는 일반 언어(입법(legislation))를 통한 전달의
확실성 사이의 구분이 이 순진한 대비가 시사하는 것만큼 확고하지 않다는
중요한 사실을 점차 인식하거나(때때로는 과장되게) 강조해온 과정이었다.
언어적으로 정식화된 일반 규칙이 사용될 때조차, 그 규칙이 요구하는 행위의
형태에 관한 불확실성은 특정한 구체적 사례에서 발생할 수 있다. 구체적인
사실 상황은 일반 규칙의 적용 여부가 문제되는 사안의 사례들로 미리 서로
구분되고 명명되어 있는 것이 아니며, 규칙 자체가 나서서 스스로의 사례라고
주장할 수도 없다. 규칙의 영역뿐 아니라 모든 경험의 영역에서는, 일반
언어가 제공할 수 있는 지침에는 언어의 본성에 내재한 한계가 존재한다.
기실 유사한 맥락에서 반복적으로 발생하며 일반적 표현이 분명히 적용가능한
명백한 사례들(예: “무엇이 탈 것(a vehicle)이라면, 자동차는 틀림없이 그
중 하나다(If anything is a vehicle a motor-car is one)”)이 존재하지만,
그 표현이 적용되는지 불분명한 사례들(예: “여기서 ‘탈 것(vehicle)’이라는
표현은 자전거(bicycles)나 비행기(airplanes), 롤러스케이트(roller
skates)도 포함하는가?”)도 있다. 후자의 경우는, 자연이나 인간의 발명에
의해 끊임없이 등장하며, 명백한 사례의 일부 특징만을 가지되 다른 특징은
결여한 사실 상황들이다. ‘해석’의 준칙들(canons of 'interpretation')은
이러한 불확실성을 줄일 수는 있어도 없앨 수는 없다. 왜냐하면 이 준칙들
역시 언어 사용에 관한 일반 규칙들이며, 그 자체가 해석을 요하는 일반
용어들을 사용하기 때문이다. 이러한 규칙들 역시 다른 규칙들과 마찬가지로
자기 자신의 해석을 제공할 수 없다. 해석이 필요 없어 보이며 사례들의
식별이 문제없거나 ‘자동적’으로 보이는 명백한 사례들은, 사실 유사한
맥락에서 반복적으로 나타나는 친숙한 것들일 뿐이다. 이러한 경우에는 분류
용어의 적용에 대해 판단이 일치하는 일반적 합의가 존재한다.

일반 용어들(general terms)은 그러한 친숙하고 일반적으로 다투어지지 않는
사례들이 없다면 의사소통의 매체로서 우리에게 무용할 것이다. 그러나
친숙한 것의 변형들(variants) 역시, \textbf{(p.~127)} 어느 특정 시점에
우리의 언어적 자원의 일부를 이루는 일반 용어들 아래에서의
분류(classification)를 요구한다. 여기서 의사소통에서 일종의 위기가
촉발된다. 일반 용어를 사용하는 데에는 찬성하는 이유들과 반대하는
이유들이 모두 있으며, 분류해야 하는 당사자에게 그 용어의 사용을
명하거나, 반대로 그 사용의 거부를 명하는 확고한 관행이나 일반적 합의는
없다. 이러한 경우에 의문들이 해소되어야 한다면, 그것들을 해소해야 할
자는 열린 대안들(open alternatives) 사이에서 일종의 선택을 해야만 한다.

이 지점에서 규칙이 표현되어 있는 권위 있는 일반 언어(authoritative
general language)는 권위 있는 예시가 그러하듯이 불확실한 방식으로만
인도할 수 있다. 이 지점에서는 규칙의 언어가 우리로 하여금 쉽게 식별
가능한 사례들을 단순히 골라낼 수 있게 해 줄 것이라는 느낌이 무너진다.
포섭(subsumption)과 삼단논법적 결론의 도출(drawing of a syllogistic
conclusion)은 무엇이 해야 할 옳은 일인지를 정하는 데 관여하는 추론의
핵심을 더 이상 특징짓지 못한다. 대신 규칙의 언어는 이제 권위 있는 예시
하나, 곧 명백한 사례(plain case)에 의해 구성되는 예시를 지목하는 것처럼
보일 뿐이다. 이것은 선례와 거의 같은 방식으로 사용될 수 있지만, 규칙의
언어는 선례보다 더 지속적이고 더 촘촘하게 주목을 요구하는 특징들을
제한한다. 공원에서 차량 사용을 금지하는 규칙이 그 적용 여부가
불확정적으로 보이는 어떤 상황들의 결합에 적용되는가라는 질문에 직면할
때, 답해야 하는 사람이 할 수 있는 전부는 (선례를 사용하는 사람이
그러하듯이) 현재 사례가 명백한 사례와 ‘관련된’ 측면들에서 ‘충분히’
유사한지를 고려하는 것이다. 언어가 그에게 남겨 두는 이러한
재량(discretion)은 매우 넓을 수 있다. 그래서 그가 그 규칙을 적용한다면,
그 결론은 자의적이거나 비합리적이지 않을 수 있음에도 불구하고 사실상
하나의 선택이다. 그는 법적으로 관련되고 충분히 가까운 것으로 합리적으로
옹호될 수 있는 유사성들 때문에, 하나의 새로운 사례를 사례들의 계열(line
of cases)에 추가하기로 선택한다. 법규칙들의 경우, 관련성(relevance)과
유사성의 가까움(closeness of resemblance)의 기준들은 법체계 전반을
관통하는 많은 복잡한 요소들과 그 규칙에 귀속될 수 있는 목적들 또는
취지(aims or purpose)에 의존한다. 이것들을 특징짓는 것은 법적 추론(legal
reasoning)에 특유하거나 고유한 것이 무엇이든 그것을 특징짓는 일이 될
것이다.

행위 기준들의 전달을 위해 선례(precedent)든 입법(legislation)이든 어느
장치가 선택되더라도, 그것들은 보통 사례들의 대다수에서는 아무리 매끄럽게
작동한다 하더라도, 그 적용이 문제되는 어떤 지점에서는
불확정적(indeterminate)임이 드러날 것이다. 그것들은 이른바 \emphksans{개방적
직조(open texture)}를 지닐 것이다. 지금까지 우리는 입법의 경우에 관하여
이것을 인간 언어의 일반적 특징으로 제시해 왔다. 경계선(borderline)에서의
불확실성은 사실의 문제들에 관한 어떤 형태의 의사소통에서든 일반적 분류
용어들을 사용하는 데 치러야 할 대가이다. 영어와 같은 자연언어는 그렇게
사용될 때 환원 불가능하게 개방적 직조를 지닌다. 그러나 실제 언어가
개방적 직조의 특성들을 가진다는 점에 의존한다는 이 사실과는 별도로, 특정
사례에 적용되는지 여부가 언제나 사전에 확정되어 있고 실제 적용 지점에서
열린 대안들 사이의 새로운 선택을 전혀 수반하지 않는, 그토록 상세한
규칙이라는 구상을---이상(ideal)으로서조차---왜 우리가 품어서는 안
되는지를 이해하는 것이 중요하다. 짧게 말해 그 이유는, 그러한 선택의
필요성이 우리가 신들(gods)이 아니라 인간들(men)이기 때문에 우리에게
떠맡겨지기 때문이다. 특정한 경우마다 추가적인 공무담당 지시(official
direction) 없이 사용될 일반 기준들에 의해, 어떤 행위 영역을 모호성 없이
사전에 규율하려 할 때마다, 우리가 서로 연결된 두 가지 불리한 조건
아래에서 애쓴다는 점은 인간적 처지(the human predicament)의 특징이며,
따라서 입법적 처지의 특징이다. 첫 번째 불리한 조건은 사실에 대한 우리의
상대적 무지(relative ignorance of fact)이고, 두 번째는 목적의 상대적
불확정성(relative indeterminacy of aim)이다. 만약 우리가 사는 세계가
유한한 수의 특징들로만 특징지어지고, 이 특징들과 그것들이 결합할 수 있는
모든 방식이 우리에게 알려져 있다면, 모든 가능성에 대해 사전에 대비할 수
있을 것이다. 우리는 특정 사례들에 대한 적용이 결코 추가적 선택을
요구하지 않는 규칙들을 만들 수 있을 것이다. 모든 것이 알려질 수 있을
것이고, 모든 것에 대해, 그것이 알려질 수 있기 때문에, 무언가가 행해질 수
있으며 규칙에 의해 사전에 특정될 수 있을 것이다. 이것은
‘기계적(mechanical)’ 법리학(jurisprudence)에 알맞은 세계일 것이다.

명백히 이 세계는 우리의 세계가 아니다. 인간 입법자들은 미래가 가져올 수
있는 모든 가능한 상황 조합들에 관해 그러한 지식을 가질 수 없다. 이러한
예견할 수 없음(inability to anticipate)은 목적의 상대적
불확정성(indeterminacy of aim)을 동반한다. 우리가 어떤 일반적 행위
규칙을 과감히 정식화할 때(예: 어떤 차량도 공원에 들어와서는 안 된다는
규칙), \textbf{(p.~129)} 이 맥락에서 사용되는 언어는 어떤 것이 그 범위
안에 들려면 충족해야 하는 필요조건들을 고정하며, 확실히 그 범위 안에
드는 것의 몇몇 명확한 예들이 우리의 마음에 떠오를 수 있다. 그것들은
전범적(paradigm) 명백한 사례들(clear cases), 곧 자동차, 버스,
오토바이이다. 그리고 우리의 입법 목적은 지금까지는
확정적(determinate)인데, 이는 우리가 일정한 선택을 했기 때문이다. 우리는
적어도 이러한 것들을 배제하는 대가를 치르더라도 공원의 평온과 조용함이
유지되어야 한다는 문제를 처음에 확정해 두었다. 다른 한편, 우리가 공원의
평온이라는 일반 목적을 처음에 예견하지 않았거나 어쩌면 예견할 수 없었던
사례들(예컨대 전기로 움직이는 장난감 자동차)과 결합시키기 전까지, 이
방향에서 우리의 목적은 불확정적(indeterminate)이다. 우리는 예견되지 않은
사례가 발생할 때 제기될 문제를, 우리가 예견하지 않았기 때문에, 아직
확정해 두지 않았다. 즉, 이러한 것들을 사용하는 것이 즐거움이나 이익이
되는 아이들을 위해 공원의 어느 정도 평온을 희생할 것인지, 아니면 그
아이들에 맞서 그 평온을 방어할 것인지의 문제이다. 예견되지 않은 사례가
실제로 발생하면, 우리는 걸려 있는 쟁점들에 직면하고, 그때 서로 경쟁하는
이익들 사이에서 우리에게 가장 만족스러운 방식으로 선택함으로써 그 문제를
확정할 수 있다. 그렇게 함으로써 우리는 우리의 최초 목적을 더 확정적으로
만들게 되며, 부수적으로는 이 규칙의 목적상 어떤 일반 단어의 의미에 관한
문제도 확정하게 된다.

서로 다른 법체계들이나 동일한 법체계라도 서로 다른 시기에는, 일반
규칙들을 특정 사례들에 적용할 때 추가적인 선택 행사가 필요하다는 점을
어느 정도로든, 그리고 더 많거나 더 적게 명시적으로 무시하거나 인정할 수
있다. 법이론에서 형식주의(formalism) 또는 개념주의(conceptualism)로
알려진 폐해는, 일반 규칙이 일단 설정되고 나면 그러한 선택의 필요성을
은폐하고 최소화하려는 언어로 정식화된 규칙들에 대한 태도에 있다. 이를
수행하는 한 가지 방식은 규칙의 의미를 동결하여, 그 일반 용어들이 적용이
문제되는 모든 경우에서 동일한 의미를 가져야 한다고 하는 것이다. 이를
확보하기 위해 우리는 명백한 사례(plain case)에 존재하는 특정 특징들을
붙잡고, 어떤 것이 그러한 특징들을 가지면 그것이 다른 어떤 특징들을 더
가지거나 결여하든, 그리고 이 방식으로 규칙을 적용할 때 사회적 결과가
무엇이든, 그것을 규칙의 범위 안에 포함시키기에 그러한 특징들이
필요충분하다고 주장할 수 있다. 이렇게 하는 것은, 우리가 그 구성에 대해
무지한 장래의 사례들의 범위에서 무엇이 행해져야 하는지를 맹목적으로 미리
판단하는 대가로 일정한 확실성 또는 \textbf{(p.~130)}
예측가능성(predictability)을 확보하는 것이다. 따라서 우리는 기실 사전에
쟁점들을 확정하는 데 성공하겠지만, 동시에 어둠 속에서, 그것들이 발생하고
식별될 때에만 합리적으로 확정될 수 있는 쟁점들을 확정하게 된다. 이러한
기법은 합리적인 사회적 목적들에 효력을 부여하기 위해 우리가 배제하고
싶어 할 사례들, 그리고 우리가 그것들을 덜 경직적으로 정의된 채로
두었더라면 우리 언어의 개방적 직조(open texture)를 지닌 용어들이 우리로
하여금 배제할 수 있게 했을 사례들을, 규칙의 범위 안에 포함하도록 우리를
강제할 것이다. 이로써 우리의 분류들의 경직성은 그 규칙을 가지거나
유지하려는 우리의 목적들과 충돌하게 될 것이다.

이 과정의 완성은 법리학자들의 ‘개념의 천국(heaven of concepts)’이다.
이는 하나의 일반 용어가 단일한 규칙의 모든 적용에서 동일한 의미를 갖는
데 그치지 않고, 법체계 내의 어떤 규칙에서 등장하든 언제나 동일한 의미를
부여받을 때에 도달한다. 그러한 경우에는, 그 용어가 여러 차례 재현될
때마다 서로 다른 이해관계나 쟁점을 고려하여 그 의미를 해석하려는 어떠한
노력도 더 이상 요구되거나 시도되지 않는다.

사실 모든 체계들은 각기 다른 방식으로 두 가지 사회적 필요 사이에서
타협한다. 하나는 넓은 행위 영역들에 걸쳐, 새로운 공무담당 지침(official
guidance)이나 사회적 쟁점들의 형량 없이도 사적 개인들이 자기 자신에게
안전하게 적용할 수 있는 확실한(certain) 규칙들에 대한 필요이고, 다른
하나는 구체적 사례에서 발생할 때에만 비로소 적절하게 이해되고 확정될 수
있는 쟁점들을, 나중에 정보를 갖춘 공무담당 선택(official choice)에 의해
확정되도록 열어 두어야 할 필요이다. 어떤 법체계들에서는 어떤 시기에는
확실성(certainty)에 지나치게 많은 것이 희생되어, 제정법들(statutes)이나
선례에 대한 사법적 해석이 지나치게 형식적이고, 그래서 사회적 목적들의
관점에서 고려될 때에만 보이는 사례들 사이의 유사성과 차이에 응답하지
못할 수 있다. 다른 체계들이나 다른 시기에는, 법원이 선례들에서 지나치게
많은 것을 영구히 열려 있거나 수정 가능한 것으로 취급하고, 입법 언어가 그
개방적 직조에도 불구하고 결국 제공하는 한계들에 지나치게 적은 존중을
보이는 것처럼 보일 수 있다. 이 문제에서 법이론은 기묘한 역사를 가진다.
그것은 법규칙들의 불확정성들(indeterminacies)을 무시하거나 과장하는
경향이 있기 때문이다. 이러한 극단들 사이의 진동에서 벗어나려면, 이
\textbf{(p.~131)} 불확정성의 뿌리에 놓인 미래를 예견하지 못하는 인간적
한계가 행위의 서로 다른 영역들에서 정도를 달리하며, 법체계들이 그 한계에
상응하는 다양한 기법들로 대응한다는 점을 상기해야 한다.

때로 법적으로 통제되어야 할 영역은 처음부터, 개별 사례들의 특징들이
사회적으로 중요한 그러나 예측 불가능한 측면들에서 너무나 다양하게 변할
것이어서, 추가적인 공무담당 지시(official direction) 없이 사례에서
사례로 적용될 일양적 규칙들을 입법부가 사전에 유용하게 정식화할 수 없는
영역으로 인정된다. 따라서 그러한 영역을 규율하기 위해 입법부는 매우
일반적인 기준들을 세우고, 다양한 사례 유형들에 정통한 행정적 규칙제정
기관에, 그 특수한 필요들에 맞추어진 규칙들을 만들어 내는 과업을
위임한다. 예컨대 입법부는 어떤 산업에 일정한 기준들을 유지하라고 요구할
수 있다. 즉 \emphksans{공정한 요율(fair rate)}만을 부과하거나 \emphksans{안전한
작업 체계들(safe systems)}을 제공하라고 요구할 수 있다. 서로 다른
기업들이 이 모호한 기준들을 스스로에게 적용하도록 맡겨 두어,
\emphksans{사후적으로(ex post facto)} 그것들을 위반한 것으로 판정될 위험을
부담하게 하기보다는, 행정기관이 규정으로 특정 산업에 대해 무엇이 ‘공정한
요율’ 또는 ‘안전한 체계’로 간주될 것인지를 특정할 때까지 위반에 대한
제재 사용을 유보하는 것이 최선이라고 판명될 수 있다. 이러한 규칙제정
권한은 해당 산업에 관한 사실들에 대한 사법적 조사와 유사한 것 및 특정한
규제 형식에 찬반하는 논변들의 청문을 거친 뒤에만 행사될 수 있을지도
모른다.

물론 매우 일반적인 기준들(very general standards)이 있더라도, 무엇이
그것들을 충족하거나 충족하지 않는지에 관한 명백하고 다툼의 여지가 없는
예들은 있을 것이다. 무엇이 ‘공정한 요율’ 또는 ‘안전한 체계’인지, 또는
아닌지에 관한 몇몇 극단적 사례들은 언제나 \emphksans{처음부터(ab initio)}
식별 가능할 것이다. 따라서 무한히 다양한 사례들의 범위 한쪽 끝에는 필수
서비스에 대해 공중을 사실상 인질로 잡으면서 기업가들에게 막대한 이윤을
안겨주는 너무 높은 요율이 있을 것이고, 다른 끝에는 기업 운영의 유인을
제공하지 못하는 너무 낮은 요율이 있을 것이다. 이 둘은 서로 다른 방식으로
요율 규제에서 우리가 가질 수 있는 어떤 가능한 목적도 좌절시킬 것이다.
그러나 이것들은 서로 다른 요소들의 범위에서 극단들일 뿐이며 실제로
마주칠 개연성이 높지 않다. 그 사이에 주의를 요구하는 어려운 실제
사례들이 놓여 있다. 관련 요소들의 예견 가능한 조합들은 적고, 이는 공정한
요율 또는 안전한 체계라는 우리의 최초 목적에 상대적 불확정성을 수반하며,
추가적인 공무담당 \textbf{(p.~132)} 선택(official choice)의 필요를
수반한다. 이러한 경우에는 규칙제정 기관(rule-making authority)이
재량(discretion)을 행사해야 한다는 점이 명확하며, 여러 사례들이 제기하는
문제를 많은 상충하는 이익들 사이의 합리적 타협인 답과는 구별되는,
발견되어야 할 하나의 유일하게 올바른 답이 있는 것처럼 다룰 가능성은
없다.

두 번째 유사한 기법은, 통제되어야 할 영역이 일률적으로 행해지거나 삼가야
할 특정 행위들의 부류를 식별하여 그것들을 단순 규칙(simple rule)의
대상으로 삼는 것은 불가능하지만, 상황들의 범위가 매우 다양하더라도 공통
경험의 친숙한 특징들을 포괄하는 경우에 사용된다. 여기서는 무엇이
‘합리적(reasonable)’인지에 관한 공통 판단들이 법에 의해 사용될 수 있다.
이 기법은 여러 예견 불가능한 형태로 발생하는 사회적 요구들(social
claims) 사이에서 형량하고 합리적 균형을 맞추는 과업을, 법원의 교정을
받을 수 있는 개인들에게 남겨 둔다. 이 경우 개인들은 가변적 기준(variable
standard)이 공식적으로 규정되기 \emphksans{전에(before)} 그 기준에 부합할
것을 요구받으며, 그들이 그것을 위반했을 때에만 \emphksans{사후적으로(ex post
facto)} 법원으로부터, 특정 행위들이나 부작위들의 관점에서 자신들에게
요구된 기준이 무엇이었는지를 알게 될 수 있다. 그러한 사안들에 대한
법원의 결정들이 선례로 간주되는 곳에서는, 가변적 기준에 대한 법원의
특정화는 행정기관이 위임된 규칙제정 권한(delegated rulemaking power)을
행사하는 것과 매우 유사하지만, 또한 명백한 차이들도 있다.

영미법(Anglo-American law)에서 이 기법의 가장 유명한 예는
과실(negligence) 사건들에서 상당한 주의(due care)의 기준을 사용하는
것이다. 타인에게 신체적 손해를 입히는 것을 피하기 위해 합리적 주의를
기울이지 않은 자들에게는 민사 제재가, 더 드물게는 형사 제재가 적용될 수
있다. 그러나 구체적 상황에서 합리적 주의 또는 상당한 주의란 무엇인가?
물론 우리는 상당한 주의의 전형적 예들을 들 수 있다. 예컨대 통행이
예상되는 곳에서 멈추고, 살피고, 듣는 것과 같은 일들이다. 그러나 우리
모두는 주의가 요구되는 상황들이 엄청나게 다양하며, 이제는 ‘멈추고,
살피고, 듣기’ 외에 또는 그 대신에 많은 다른 행위들이 요구된다는 점을 잘
알고 있다. 기실 보는 것이 위험을 피하는 데 도움이 되지 않는다면 이것들은
충분하지 않을 수 있고 전혀 쓸모없을 수도 있다. 합리적 주의 기준들을
적용하면서 우리가 추구하는 것은 (1) 중대한 해악을 피하게 할
\textbf{(p.~133)} 예방조치들이 취해지게 하면서도, (2) 그 예방조치들이
적절한 예방조치들의 부담이 다른 존중할 만한 이익들의 지나치게 큰 희생을
수반하지 않도록 하는 것이다. 물론 피를 흘리며 죽어가는 사람이 병원으로
이송되는 경우가 아니라면, 멈추고, 살피고, 듣는 것에 의해 그리 많은 것이
희생되지는 않는다. 그러나 주의가 요구되는 가능한 사례들의 엄청난 다양성
때문에, 해악에 대한 예방조치가 취해져야 한다면 어떤 상황들의 조합이
발생할지, 또는 어떤 이익들이 어느 정도로 희생되어야 할지를 우리는
\emphksans{처음부터(ab initio)} 예견할 수 없다. 따라서 우리는 특정 사례들이
발생하기 전에는, 해악의 위험을 줄이기 위해 이익들이나 가치들의 어떤 희생
또는 타협을 정확히 하고자 하는지를 고려할 수 없다. 다시 말해, 사람들을
해악으로부터 보호하려는 우리의 목적은, 경험만이 우리 앞에 가져다줄
가능성들과 결합시키거나 그것들에 비추어 시험하기 전까지는
불확정적(indeterminate)이다. 경험이 그것을 가져오면, 그때 우리는 하나의
결정에 직면해야 하며, 그 결정은 내려질 때 우리의 목적을 \emphksans{그만큼(pro
tanto)} 확정적으로 만들 것이다.

이 두 가지 기법(techniques)을 고려하면, 가변적 기준(variable standard)이
아니라 특정 행위들을 요구하는 규칙(rule)에 의해, 개방적 직조(open
texture)의 가장자리(fringe)만을 남긴 채 \emphksans{처음부터(ab initio)}
성공적으로 통제되는 넓은 행위 영역들의 특징들이 두드러진다. 그것들은
다음 사실에 의해 특징지어진다. 어떤 구별 가능한 행위들, 사건들, 또는
사태들이 우리에게 회피되어야 하거나 생겨나야 할 것으로서 실천적으로 매우
중요하여, 그것들과 함께 나타나는 부수적 상황들이 우리로 하여금 그것들을
달리 보도록 기울게 하는 경우가 거의 없다는 것이다. 이것의 가장 거친 예는
인간 존재의 살해(killing of a human being)이다. 인간 존재들이 다른 인간
존재들을 살해하는 상황들은 매우 다양하지만, 우리는 가변적 기준(‘인간
생명에 대한 상당한 존중’)을 세우는 대신 살해를 금지하는 규칙을 만들 수
있는 위치에 있다. 이는 생명 보호의 중요성에 대한 우리의 평가보다 더 큰
무게를 갖거나 그 평가를 수정하게 하는 요소들이 우리에게 거의 없는 것으로
보이기 때문이다. 거의 언제나 살해는 말하자면 그것에 수반되는 다른
요소들을 \emphksans{지배한다(dominates)}. 따라서 우리가 그것을
‘살해(killing)’로서 사전에 배제할 때, 우리는 서로 형량되어야 하는
쟁점들을 맹목적으로 미리 판단하고 있는 것이 아니다. 물론 예외들, 이
통상적으로 지배적인 요소를 압도하는 요소들이 있다. 정당방위에서의 살해와
기타 정당화 가능한 살해(justifiable homicide)의 형태들이 그것이다.
그러나 이것들은 적고 상대적으로 단순한 용어들로 식별 가능하며, 일반
규칙에 대한 예외들로 인정된다.

\textbf{(p.~134)} 쉽게 식별 가능한 어떤 행위, 사건, 또는 사태의 지배적
지위가 어떤 의미에서는 관행적(conventional)이거나 인위적(artificial)일
수 있으며, 인간 존재로서 우리에게 그것이 가지는 ‘자연적(natural)’ 또는
‘내재적(intrinsic)’ 중요성 때문이 아닐 수 있다는 점을 주목하는 것이
중요하다. 도로 통행 규칙(rule of the road)이 어느 쪽 도로를 지정하는지는
중요하지 않으며, 부동산 양도증서(conveyance)의 작성(execution)을 위해
어떤 형식요건들이 규정되는지도 (일정 한계 안에서는) 중요하지 않다.
그러나 이러한 사항들에 대해 쉽게 식별 가능하고 일양적인 절차가 있어야
하며, 따라서 명확한 옳고 그름(right and wrong)이 있어야 한다는 점은
대단히 중요하다. 이것이 법에 의해 도입되면, 소수의 예외를 제외하고는
그것을 고수하는 것의 중요성이 최우선적(paramount)이 된다. 왜냐하면
상대적으로 적은 수의 부수적 상황들만이 그것보다 더 큰 무게를 가질 수
있고, 그러한 상황들은 쉽게 예외로 식별되어 규칙으로 환원될 수 있기
때문이다. 영국 부동산법(real property law)은 규칙들의 이러한 측면을 매우
명확하게 예증한다.

권위 있는 예시들(authoritative examples)에 의한 일반 규칙들의 전달은,
우리가 보았듯이, 더 복잡한 종류의 불확정성들을 수반한다. 선례를 법적
유효성의 기준으로 인정하는 것은 서로 다른 체계들에서, 그리고 동일한
체계의 서로 다른 시기들에서 서로 다른 것들을 의미한다. 영국의 선례
‘이론(theory)’에 대한 서술들은 어떤 지점들에서는 여전히 매우 논쟁적이다.
기실 그 이론에서 사용되는 핵심 용어들, 즉 ‘\emphksans{판결 이유(ratio
decidendi)}’, ‘중요 사실(material facts)’, ‘해석(interpretation)’조차도
각자 고유한 불확실성의 반음영(penumbra)을 지닌다. 우리는 어떤 새로운
일반적 서술을 제시하지는 않고, 제정법(statute)의 경우에서와 마찬가지로,
개방적 직조(open texture)의 영역과 그 안에서 이루어지는 창조적 사법
활동을 간략하게 특징짓고자 할 뿐이다.

영국법에서 선례(precedent)의 사용에 대한 정직한 서술은 다음과 같은
대조적 사실들의 쌍들을 위한 자리를 반드시 마련해야 한다. \emphksans{첫째},
주어진 권위 있는 선례(authoritative precedent)가 어떤 규칙에 대한
권위(authority)인지를 결정하는 단일한 방법은 없다. 그럼에도 불구하고
판결이 내려진 사건들의 압도적 다수에서는 의문이 거의 없다.
사건요지(head-note)는 보통 충분히 정확하다. \emphksans{둘째}, 사건들로부터
추출될 어떤 규칙에 대해서도 권위 있거나 유일하게 올바른
정식화(authoritative or uniquely correct formulation)는 없다. 다른 한편,
선례가 후속 사건에 대해 가지는 관련성(bearing)이 문제될 때, 주어진
정식화가 적절하다는 데 매우 일반적인 합의가 있는 경우가 흔하다.
\emphksans{셋째}, 선례에서 추출된 규칙이 어떤 권위 있는 지위(authoritative
status)를 가지든, 그것은 그 규칙에 구속되는 법원들이 \textbf{(p.~135)}
다음 두 유형의 창조적 또는 입법적 활동을 행사하는 것과 양립 가능하다.
한편으로 후속 사건을 판단하는 법원들은 선례에서 추출된 규칙을 좁히고,
이전에 고려되지 않았거나 고려되었더라도 열려 있던 어떤 예외를
인정함으로써 선례와 반대되는 결정에 이를 수 있다. 이처럼 앞선 사건을
‘구별(distinguishing)’하는 과정은 그 사건과 현재 사건 사이에서 법적으로
관련된 어떤 차이를 발견하는 것을 포함하며, 그러한 차이들의 부류는 결코
소진적으로 확정될 수 없다. 다른 한편, 법원들은 앞선 선례를 따르면서도,
앞선 사건에서 정식화된 규칙 안에서 발견되는 어떤 제약(restriction)을,
그것이 제정법(statute)이나 더 앞선 선례에 의해 확립된 어떤 규칙에
의해서도 요구되지 않는다는 근거로 폐기할 수 있다. 이렇게 하는 것이
규칙을 넓히는 것이다. 선례의 구속력(binding force)에 의해 열려 있는
이러한 두 형태의 입법적 활동에도 불구하고, 영국 선례 체계의 결과는 그
사용을 통해 중요도가 큰 것과 작은 것을 모두 포함하는 매우 많은 수의
규칙들의 집합을 산출한 것이며, 그 가운데 많은 규칙들은 어떤 제정법
규칙(statutory rule) 못지않게 확정적(determinate)이다. 이제 그것들은
제정법(statute)에 의해서만 변경될 수 있으며, 법원들 자신도 확립된
선례들의 요구와 ‘실질적 타당성(merits)’이 배치되는 것처럼 보이는
사건들에서 흔히 그렇게 선언한다.

법의 개방적 직조(open texture)는, 기실 많은 행위 영역들에서 사안마다
무게가 달라지는 상충하는 이익들 사이의 균형을 상황에 비추어 맞추는
법원들이나 공무담당자들에 의해 많은 것이 발전되도록 남겨져야 함을
의미한다. 그럼에도 불구하고 법의 삶(life of the law)은 매우 큰 정도로,
가변적 기준들(variable standards)의 적용과는 달리 사례마다 새로운 판단을
요구하지 않는 확정적 규칙들(determinate rules)에 의한 공무담당자들과
사적 개인들 모두의 지도로 이루어진다. 이러한 사회생활의 두드러진 사실은,
어떤 규칙이든(문서화된 것이든 선례에 의해 전달된 것이든) 구체적 사례에
대한 적용가능성을 둘러싸고 불확실성들이 터져 나올 수 있음에도 불구하고
여전히 참이다. 여기, 규칙들의 주변부에서 그리고 선례 이론에 의해 열려
있는 영역들에서, 법원들은 행정기관들이 가변적 기준들을 정교화할 때
중심적으로 수행하는 규칙-생산 기능을 수행한다. \emphksans{선례구속 원칙(stare
decisis)}이 확고히 인정된 체계에서, 법원의 이 기능은 행정기관에 의한
위임된 규칙제정 권한(delegated rule-making powers)의 행사와 매우
유사하다. 영국에서는 이 사실이 형식들(forms)에 의해 자주 가려진다.
법원들이 \textbf{(p.~136)} 흔히 그러한 창조적 기능을 부인하고, 제정법
해석과 선례 사용의 적절한 과업은 각각 ‘입법자의 의도(intention of the
legislature)’와 이미 존재하는 법을 탐색하는 것이라고 고수하기 때문이다.

\subsection{2. 규칙회의주의(rule-scepticism)의
다양상들}\label{uxaddcuxce59uxd68cuxc758uxc8fcuxc758rule-scepticismuxc758-uxb2e4uxc591uxc0c1uxb4e4}

우리는 법의 개방적 직조(open texture)에 대해 상당히 길게 논의했는데,
이는 이 특징을 정당한 관점에서 보는 것이 중요하기 때문이다. 그것을
정당하게 다루지 못하면(do justice) 언제나 법의 다른 특징들을 흐리게 할
과장들이 유발될 것이다. 모든 법체계에는 처음에는 모호한 기준들을
확정적으로 만들고, 제정법들(statutes)의 불확실성들을 해소하며, 권위 있는
선례들(authoritative precedents)에 의해 단지 넓게만 전달된 규칙들을
발전시키고 한정하는 데서, 법원들과 기타 공무담당자들이
재량(discretion)을 행사하도록 넓고 중요한 영역이 열려 있다. 그럼에도
불구하고 이러한 활동들은, 중요하고 충분히 연구되지 않았더라도, 그것들이
일어나는 틀(framework)과 그것들의 주된 최종 산물이 모두 일반 규칙들로
이루어져 있다는 사실을 가려서는 안 된다. 이것들은 개인들이 사례마다 그
적용을 스스로 볼 수 있는 규칙들로서, 추가적인 공무담당 지시(official
direction)나 재량에 더 호소할 필요가 없다.

규칙들이 법체계의 구조에서 중심적 자리를 가진다는 논지(contention)가
진지하게 의심받을 수 있었다는 것은 이상해 보일 수 있다. 그러나
‘규칙회의주의(rule-scepticism)’, 곧 규칙에 관한 말은 법이 단지 법원의
결정들과 그것들에 대한 예측들로 이루어진다는 진실을 가리는(cloaking)
하나의 신화라는 주장(claim)은 법률가의 솔직함에 강력하게 호소할 수 있다.
이 주장이 이차 규칙들과 일차 규칙들을 모두 포괄하도록 무제한적인 일반
형식으로 진술된다면, 그것은 기실 매우 비정합적(incoherent)이다. 법원의
결정들이 존재한다는 단언은 아무 규칙도 전혀 존재하지 않는다는 부정과
일관되게 결합될 수 없기 때문이다. 이는 우리가 보았듯이 법원의 존재가
관할권(jurisdiction)을 계속 바뀌는 일련의 개인들에게 부여하고 그래서
그들의 결정들을 권위 있는 것으로 만드는 이차 규칙들의 존재를 수반하기
때문이다. 어떤 사람들의 공동체가 결정과 결정에 대한 예측이라는
관념들(notions)은 이해하지만 규칙이라는 관념(notion)은 이해하지
못한다면, 그들에게는 \emphksans{권위 있는(authoritative)} 결정이라는
관념(idea)이 결여될 것이며 그와 함께 법원이라는 관념(idea)도 결여될
것이다. 사인의 결정과 법원의 결정을 구별할 아무것도 없을 것이다. 우리는
\textbf{(p.~137)} ‘습관적 복종’이라는 관념(notion)으로, 법원에서
요구되는 권위 있는 관할권의 토대로서 결정의 예측가능성이 가지는 결함을
보충하려 할지도 모른다. 그러나 이렇게 하면, 입법 권한들을 부여하는
규칙의 대체물로 제4장에서 습관을 검토했을 때 드러난 모든 부적합성들을,
이 목적에서도, 습관이라는 관념(notion)이 가진다는 것을 발견하게 될
것이다.

이 이론의 더 온건한 몇몇 형태에서는, 법원들이 존재하려면 그것들을
구성하는 법규칙들이 있어야 하며, 따라서 이러한 규칙들 자체는 단순히
법원의 결정들에 대한 예측들일 수 없다는 점이 인정될 수 있다. 그러나 이
양보만으로는 실제로 거의 진전을 이룰 수 없다. 이러한 유형의 이론의
특징적 단언은 제정법들(statutes)이 법원에 의해 적용되기 전에는 법이
아니라 단지 법의 원천들(sources of law)이라는 것이고, 이는 존재하는
유일한 규칙들은 법원들을 구성하는 데 필요한 규칙들이라는 단언과 일관되지
않기 때문이다. 계속 바뀌는 일련의 개인들에게 입법 권한들을 부여하는 이차
규칙들도 있어야 한다. 왜냐하면 이 이론은 제정법들(statutes)이 있다는
것을 부정하지 않기 때문이다. 기실 그것은 그것들을 법의 단순한
‘원천들(sources)’로 인용하고, 제정법들이 법원에 의해 적용되기 전에는
법이라는 것만 부정한다.

이러한 반론들은 중요하며, 이 이론의 경솔한 형태에 대해서는 적절하지만,
그 모든 형태에 적용되지는 않는다. 규칙회의주의는 사법적 또는 입법적
권한을 부여하는 이차 규칙들의 존재를 부정하려는 의도가 전혀 없었고,
그것들이 결정들이나 결정들에 대한 예측들에 불과하다고 보일 수 있다는
주장에 결부된 적도 전혀 없었을 수 있다. 분명히 이러한 유형의 이론이 가장
자주 의존해 온 예들은 사인들에게 책무를 부과하거나 권리들 또는 권한들을
부여하는 규칙들에서 끌어온 것이다. 그러나 규칙들이 있다는 부정과
규칙들이라 불리는 것이 법원의 결정들에 대한 예측들일 뿐이라는 단언을
이러한 방식으로 제한한다고 상정하더라도, 적어도 한 의미에서는 그것이
명백히 거짓이다. 왜냐하면 적어도 근대 국가의 어떤 행위 영역들과
관련해서는, 개인들이 우리가 내부 관점(internal point of view)이라 불렀던
행위와 태도의 전 범위를 실제로 보인다는 점은 의심될 수 없기 때문이다.
법들은 그들의 삶에서 단순히 습관들이나 법원들의 결정들 또는 다른
공무담당자들의 행위들을 예측하는 토대로 기능하는 것이 아니라, 수용된
법적 행위 기준들(accepted legal standards of behaviour)로 기능한다. 즉
그들은 단지 법이 자신들에게 요구하는 것을 \textbf{(p.~138)} 상당한
정도의 반복적 일양성(tolerable regularity)으로 행할 뿐만 아니라, 그것을
법적 행위 기준으로 바라보고, 타인들을 비판하거나 요구들을 정당화할 때
그것을 참조하며, 타인들이 자신들에게 가하는 비판들과 요구들을 받아들일
때에도 그것을 참조한다. 법규칙들을 이러한 규범적 방식으로 사용할 때
그들은 물론 법원들과 기타 공무담당자들이 해당 체계의 규칙들에 따라
반복적으로 일양적인, 따라서 예측 가능한 일정한 방식들로 계속 결정하고
행동할 것이라고 상정한다. 그러나 개인들이 법원들의 결정들이나 제재가
적용될 개연성을 기록하고 예측하는 외부 관점(external point of view)에
자신들을 한정하지 않는다는 것은 분명히 관찰 가능한 사회생활의 사실이다.
그 대신 그들은 행위의 지침으로서 법에 대한 공유된 수용을 규범적 용어들로
계속 표현한다. 우리는 제3장에서 ‘의무(obligation)’와 같은 규범적
용어들이 공무담당 행위에 대한 예측 이상의 아무것도 의미하지 않는다는
주장을 상세히 검토했다. 우리가 논변했듯이 그 주장이 거짓이라면,
법규칙들은 사회생활 속에서 그러한 것으로 기능한다. 그것들은 습관들이나
예측들의 서술이 아니라 규칙들로서 \emphksans{사용된다(used)}. 물론 그것들은
개방적 직조를 지닌 규칙들이며, 그 직조가 열린 지점들에서 개인들은
법원들이 어떻게 결정할지를 예측하고 그에 따라 자기 행위를 조정할 수 있을
뿐이다.

규칙회의주의는 우리의 진지한 주의를 요구하지만, 오직 사법적
결정(judicial decision)에서 규칙들이 수행하는 기능에 관한 이론으로서만
그렇다. 이 형태에서, 앞서 우리가 주의를 환기한 모든 반론들을
인정하면서도, 그것은 법원들에 관한 한 개방적 직조의 영역을 경계 지을
아무것도 없다는 논지(contention)에 이른다. 따라서 판사들을 그들이 하는
방식으로 사건들을 결정하도록 스스로 규칙들에 종속되어 있거나 규칙들에
‘구속되어(bound)’ 있다고 보는 것은, 무의미하지 않더라도 거짓이라는
것이다. 판사들은 충분히 예측 가능한 정도의 반복적 일양성(regularity)과
일양성(uniformity)을 가지고 행동하여, 다른 사람들이 긴 기간에 걸쳐
법원의 결정들에 의해 규칙들처럼 살아갈 수 있게 할 수 있다. 판사들은
심지어 자신들이 그렇게 결정할 때 강제감(feelings of compulsion)을 경험할
수도 있고, 이러한 감정들 역시 예측 가능할 수 있다. 그러나 이 너머에는
그들이 준수하는 규칙으로 특징지을 수 있는 것은 아무것도 없다. 법원들이
올바른 사법적 행위의 기준들(standards of correct judicial behaviour)로
취급하는 것은 아무것도 없으며, 따라서 그 행위 안에는 규칙들의 수용에
특징적인 내부 관점(internal point of view)을 드러내는 것도 아무것도
없다는 것이다.

이 이론의 이러한 형태는 매우 다양한 성격과 중요도를 지닌 근거들로부터
지지를 얻는다. 규칙회의주의자는 \textbf{(p.~139)} 때로는 실망한
절대주의자(absolutist)이다. 그는 규칙들이 형식주의자의
이상세계(formalist's heaven)에서 그러할 모습 그대로도 아니며, 인간이
신들처럼 되어 모든 가능한 사실의 조합을 예견할 수 있고 그래서 개방적
직조가 규칙들의 필수적 특징이 아닌 세계에서 그러할 모습 그대로도
아니라는 점을 발견한 것이다. 규칙이 존재한다는 것이 무엇인지에 관한
회의주의자의 구상은 도달 불가능한 이상(ideal)일 수 있고, 소위 규칙이라
불리는 것에서 그 이상이 달성되지 않는다는 것을 발견할 때, 그는 규칙들이
존재한다거나 존재할 수 있다는 것을 부정함으로써 자신의 실망을 표현한다.
이렇게 해서, 판사들이 사건을 판결할 때 자신들을 구속한다고 주장하는
규칙들이 개방적 직조를 갖고 있다거나, 사전에 소진적으로 특정될 수 없는
예외들을 가진다는 사실, 그리고 규칙으로부터의 일탈이 판사에게 물리적
제재(physical sanction)를 불러오지 않는다는 사실은 회의주의자의 주장을
확립하는 데 자주 이용된다. 이 사실들은 다음과 같은 점을 보이기 위해
강조된다. “규칙들이 중요한 것은 판사들이 무엇을 할지를 예측하는 데
도움을 주는 한에서이다. 예쁜 장난감(pretty playthings)으로서의 중요성을
제외하면 그것이 규칙들의 중요성의 전부이다.” \footnote{Llewellyn,
  \emph{The Bramble Bush} (2nd edn.), p.~9.}

이런 식으로 논변하는 것은 현실 생활의 어떤 영역에서든 규칙들이 실제로
무엇인지를 무시하는 것이다. 그것은 우리가 다음의 딜레마에 직면해 있다고
시사한다. “규칙들은 형식주의자의 천국(formalist's heaven)에서 그러할
모습 그대로이며 족쇄가 구속하듯 구속하거나, 아니면 규칙들은 없고 예측
가능한 결정들이나 행위 양식들만 있다.” 그러나 분명 이것은 거짓
딜레마다. 우리는 다음 날 친구를 방문하겠다고 약속한다. 그 날이 오자 그
약속을 지키는 것이 위험하게 아픈 누군가를 방치하는 것을 수반하게 된다.
이것이 약속을 지키지 않는 적절한 이유로 수용된다는 사실은, 약속을 지킬
것을 요구하는 규칙은 없고 약속을 지키는 데 어떤 반복적 일양성만 있을
뿐이라는 것을 결코 의미하지 않는다. 그러한 규칙들에 소진적으로 진술될 수
없는 예외들이 있다는 사실로부터, 모든 상황에서 우리가 우리의 재량에
맡겨져 있고 결코 약속을 지키도록 구속되지 않는다는 결론은 따르지 않는다.
‘\ldots 하지 않는 한(unless \ldots)’이라는 말로 끝나는 규칙도 여전히
규칙이다.

때때로 법원들을 구속하는 규칙들의 존재는, 어떤 사람이 특정한 방식으로
행위함으로써 자신에게 그렇게 행위할 것을 요구하는 규칙에 대한 수용을
드러냈는지 여부라는 문제가, 그 사람이 행위하기 전이나 행위하면서 거친
사고 과정들에 관한 심리학적 문제들과 \textbf{(p.~140)} 혼동되기 때문에
부정된다. 매우 자주, 어떤 사람이 규칙을 구속력 있는 것으로, 그리고
자신과 타인들이 자유롭게 바꿀 수 없는 것으로 수용할 때, 그는 주어진
상황에서 그것이 무엇을 요구하는지를 아주 직관적으로 보고, 먼저 그 규칙과
그것이 요구하는 것을 생각하지 않고도 그것을 행할 수 있다. 우리가
체스에서 규칙들에 부합하여 말을 움직이거나 신호등이 빨간색일 때 멈출 때,
우리의 규칙-준수(rule-complying) 행위는 종종 규칙들의 관점에서 계산에
의해 매개되지 않은, 상황에 대한 직접 반응이다. 그러한 행위들이 규칙의
진정한 적용이라는 증거는 그것들이 특정한 상황들 속에 놓여 있다는 점이다.
이 상황들 가운데 일부는 특정 행위에 앞서 있고, 다른 일부는 그 뒤를
따른다. 또 그 일부는 일반적이고 가설적인 용어들로만 진술될 수 있다.
행위하면서 우리가 규칙을 적용했다는 것을 보여 주는 이 요소들 가운데 가장
중요한 것은, \emphksans{만약(if)} 우리의 행위가 문제 제기된다면 우리는 그
행위를 그 규칙을 참조하여 정당화하려는 성향이 있다는 점이다. 그리고
규칙에 대한 우리의 수용의 진정성은 그것에 대한 우리의 과거 및 이후의
일반적 인정(acknowledgements)과 그것에의 부합(conformity)뿐 아니라, 우리
자신의 일탈과 타인의 일탈에 대한 비판에서도 드러날 수 있다. 이러한
증거나 유사한 증거에 근거하여 우리는 기실, 규칙에 대한 우리의 ‘생각
없는’ 준수(compliance)에 앞서 무엇이 해야 할 옳은 일이고 왜 그런지를
말하라는 질문을 받았다면, 정직할 경우 우리는 답변에서 그 규칙을 인용했을
것이라고 결론 내릴 수 있다. 진정으로 규칙의 준수(observance)인 행위를,
단지 우연히 그것과 일치하는(coincide with) 행위와 구별하는 데 필요한
것은, 우리의 행위가 그러한 상황들 속에 놓여 있다는 점이지, 규칙에 대한
명시적 사고가 그 행위에 수반되었다는 점이 아니다. 바로 이렇게 우리는,
수용된 규칙의 준수(compliance)로서의 성인 체스 플레이어의 수와, 단지
말을 옳은 자리로 밀어 넣은 아기의 행위를 구별할 것이다.

이것은 가장이나 ‘겉치레(window dressing)’가 불가능하다는 말이 아니며,
그것은 때때로 성공할 수도 있다. 어떤 사람이 자신이 규칙에 근거하여
행위했다고 단지 \emphksans{사후적으로(ex post facto)} 가장했는지를 가리는
검사들은, 모든 경험적 검사들처럼 본질적으로 오류 가능하지만, 고질적으로
그런 것은 아니다. 어떤 사회에서 판사들이 언제나 먼저 직관적으로 또는
‘직감으로(by hunches)’ 결정을 내리고, 그런 다음 법규칙들의 목록에서 당면
사건과 닮았다고 가장한 규칙 하나를 단지 선택할 수도 있다. 그러고 나서
그들은 이것이 자신들의 결정을 요구하는 것으로 자신들이 간주한 규칙이라고
주장할 수 있지만, 그들의 행위나 말의 다른 어떤 것도 \textbf{(p.~141)}
그들이 그것을 자신들을 구속하는 규칙으로 간주했다는 것을 시사하지 않을
수 있다. 일부 사법적 결정들은 이와 같을 수 있다. 그러나 대부분의 경우
결정들은 체스 플레이어의 수처럼, 결정의 인도 기준들로 의식적으로
받아들인 규칙들에 부합하려는 진정한 노력에 의해 도달되거나, 만약
직관적으로 도달되었다면, 판사가 사전에 준수하려는 성향을 가지고 있었고
당면 사건과의 관련성이 일반적으로 인정될 규칙들에 의해 정당화된다는 점은
분명하다.

규칙회의주의(rule-scepticism)의 마지막이자 가장 흥미로운 형태는
법규칙들의 개방적 성격이나 많은 결정들의 직관적 성격에 근거하지 않고,
법원의 결정이 권위 있는 것(something authoritative)으로서 독특한 지위를
가지며, 최고 법원들(supreme tribunals)의 경우에는 최종적(final)이라는
사실에 근거한다. 우리가 다음 절을 할애할 이 형태의 이론은,
그레이(Gray)가 \emph{The Nature and Sources of Law}에서 그렇게 자주
되풀이한 호들리(Hoadly) 주교의 유명한 구절에 암묵적으로 들어 있다.
“글로 쓰였든 말로 발화되었든 어떤 법들을 해석할 절대적 권위(absolute
authority)를 가진 자가 있다면, 사실상 모든 면에서 입법자(lawgiver)는
바로 그 사람이지 그것들을 처음 쓰거나 말한 사람이 아니다.”

\subsection{3. 사법적 결정에서의 최종성(Finality)과
무오류성(Infallibility)}\label{uxc0acuxbc95uxc801-uxacb0uxc815uxc5d0uxc11cuxc758-uxcd5cuxc885uxc131finalityuxacfc-uxbb34uxc624uxb958uxc131infallibility}

최고 법원(supreme tribunal)은 법이 무엇인지를 말하는 데 최후의 말(last
word)을 가지며, 일단 그렇게 말하면 그 법원이 ‘그르다(wrong)’는 진술은 그
체계 안에서는 아무런 귀결도 갖지 않는다. 그로 인해 누구의 권리들이나
책무들도 변경되지 않는다. 그 결정은 물론 입법(legislation)에 의해 법적
효과(legal effect)를 박탈당할 수 있지만, 이에 호소하는 것이 필요하다는
바로 그 사실은, 법에 관한 한 법원의 결정이 그르다는 진술의 공허한 성격을
보여 준다. 이러한 사실들을 고려하면, 최고 법원의 결정들의 경우 그
최종성(finality)과 무오류성(infallibility)을 구분하는 것은 현학적인
것처럼 보인다. 이것은 법원들이 결정할 때 언제나 규칙들에 구속된다는 것을
부정하는 또 다른 형태로 이어진다. “법(또는 헌법)은 법원들이 그것이라고
말하는 것이다.”

이러한 형태의 이론에서 가장 흥미롭고 교훈적인 특징은 ‘법(또는 헌법)은
법원들이 그것이라고 말하는 것이다’와 같은 진술들의 애매성(ambiguity)을
이용한다는 점, 그리고 이 이론이 일관되려면 \textbf{(p.~142)} 법에 대한
비공식 진술들과 법원의 공식 진술들 사이의 관계에 관해 제시해야 하는
설명이다. 이 애매성을 이해하기 위해 우리는 잠시 돌아서 게임의 경우에
있는 유비를 검토할 것이다. 많은 경쟁적 게임들은 공식 득점기록자(official
scorer) 없이 행해진다. 경쟁하는 이익들에도 불구하고, 선수들은 득점
규칙을 특정 사례들에 적용하는 데 상당히 잘 성공한다. 그들은 보통 자기
판단들에서 일치하며, 해소되지 않은 다툼들은 적을 수 있다. 공식
득점기록자 제도가 도입되기 전에, 어떤 선수가 한 득점 진술은, 그가
정직하다면, 그 게임에서 수용된 특정 득점 규칙을 참조하여 게임의 진행을
평가하려는 노력이다. 그러한 득점 진술들은 득점 규칙을 적용하는 내부
진술들(internal statements)이며, 선수들이 일반적으로 규칙들을 지키고 그
위반에 이의를 제기할 것임을 전제하지만, 이러한 사실들의 진술들이나
예측들은 아니다.

관습 체제에서 성숙한 법체계로의 변화와 마찬가지로, 그 판정들이 최종적인
득점기록자 제도를 마련하는 이차 규칙들(secondary rules)이 게임에
추가되면, 체계 안에는 새로운 종류의 내부 진술이 들어온다. 선수들의 득점
진술들과 달리, 득점기록자의 결정들(determinations)은 이차 규칙들에 의해
다툴 수 없는 지위를 부여받기 때문이다. 게임의 목적상 ‘득점은
득점기록자가 말하는 것이다’라는 말이 참인 것은 \emphksans{이런(this)}
의미에서이다. 그러나 득점 \emphksans{규칙}(scoring \emph{rule})은 이전 그대로
남아 있으며, 득점기록자의 책무(duty)는 그것을 자신이 할 수 있는 최선으로
적용하는 것임을 보는 것이 중요하다. ‘득점은 득점기록자가 말하는
것이다’가, 득점기록자가 자기 재량으로 적용하기로 선택한 것 외에는 득점
규칙이 없다는 뜻이라면 거짓일 것이다. 기실 그러한 규칙을 가진 게임도
있을 수 있고, 득점기록자의 재량이 어느 정도 반복적으로 일양적인 행사
양상을 보인다면 그것을 하는 데서 얼마간 즐거움을 찾을 수도 있을 것이다.
그러나 그것은 다른 게임일 것이다. 우리는 그러한 게임을 ‘득점기록자의
재량(scorer’s discretion)' 게임이라고 부를 수 있다.

득점기록자가 가져오는 분쟁의 신속하고 최종적인 처리라는 이점들이 대가를
치르고 사들여진다는 점은 명백하다. 득점기록자 제도는 선수들을
곤경(predicament)에 직면시킬 수 있다. 이전과 같이 게임이 득점 규칙에
의해 규율되기를 바라는 소망과, 그 적용이 의심스러운 곳에서 그 적용에
관하여 최종적이고 권위 있는 결정들을 바라는 소망은 충돌하는 목적들로
드러날 수 있다. 득점기록자는 정직한 \textbf{(p.~143)} 실수를 할 수도
있고, 술에 취했을 수도 있으며, 자기 능력의 최선으로 득점 규칙을 적용해야
한다는 책무(duty)를 함부로 위반할 수도 있다. 그는 이 이유들 중 어느 하나
때문에, 타자가 전혀 움직이지 않았는데도 ‘득점(run)’을 기록할 수 있다.
상급 권위기관에 대한 항소로 그의 판정들을 교정하는 장치가 마련될 수
있다. 그러나 이것은 어딘가에서 최종적이고 권위 있는 판단(judgment)으로
끝나야 하며, 그것은 오류 가능성이 있는 인간 존재들에 의해 내려질
것이므로 정직한 실수, 남용, 또는 위반의 동일한 위험을 수반할 것이다.
모든 규칙의 위반을 교정하기 위한 규칙을 마련하는 것은 불가능하다.

규칙들을 최종적이고 권위 있게 적용할 권위기관(authority)을 세우는 데
내재한 위험은 어느 영역에서든 현실화될 수 있다. 게임이라는 소박한
영역에서 현실화될 수 있는 위험들은 고려할 가치가 있는데, 그것들이 이
형태의 권위(authority)가 사용되는 곳이면 어디서나 그것을 이해하는 데
필요한 몇몇 구별들을 규칙회의주의자가 무시한다는 점을 특히 분명한
방식으로 보여 주기 때문이다. 공식 득점기록자(official scorer)가 세워지고
그의 득점 판정들(determinations)이 최종적인 것으로 만들어지면,
선수들이나 기타 비공식 인물들(non-officials)이 한 득점에 관한 진술들은
게임 안에서 아무런 지위를 갖지 않는다. 그것들은 그 결과에 무관하다.
그것들이 득점기록자의 진술과 우연히 일치하면 문제될 것이 없지만,
충돌한다면 결과를 산정할 때 무시되어야 한다. 그러나 이러한 매우 명백한
사실들은 선수들의 진술들이 득점기록자의 판정들에 대한 예측들로
분류된다면 왜곡될 것이며, 이러한 진술들이 득점기록자의 판정들과 충돌할
때 무시되는 것을 그것들이 거짓으로 드러난 판정들에 대한 예측들이었다고
말함으로써 설명하는 것은 터무니없을 것이다. 공식 득점기록자가 도입된
뒤에도, 선수는 자기 자신의 득점 진술들을 하면서 이전에 하던 일을 하고
있다. 즉 득점 규칙을 참조하여, 자신이 할 수 있는 최선으로 게임의 진행을
평가하고 있다. 득점기록자 자신도 자기 지위의 책무들을 이행하는 한 바로
이것을 하고 있다. 양자의 차이는 한쪽이 다른 쪽이 무엇을 말할지를
예측한다는 데 있지 않고, 선수들의 진술들은 득점 규칙의 비공식
적용들(unofficial applications)이어서 결과를 산정하는 데 의미가 없지만,
득점기록자의 진술들은 권위 있고 최종적이라는 데 있다. 만약 행해지는
게임이 ‘득점기록자의 재량(scorer’s discretion)'이라면, 비공식 진술들과
공식 진술들 사이의 관계가 \textbf{(p.~144)} 필연적으로 달라진다는 점을
관찰하는 것이 중요하다. 선수들의 진술들은 득점기록자의 판정들에 대한
예측일 \emphksans{것(would)}일 뿐 아니라, 그 밖의 아무것도 \emphksans{될 수
없을(could)} 것이다. 왜냐하면 그 경우 ‘득점은 득점기록자가 말하는
것이다’ 자체가 득점 \emphksans{규칙}일 것이기 때문이다. 선수들의 진술들이
득점기록자가 공식적으로 하는 것의 단순한 비공식 판본들일 가능성은 없을
것이다. 그러면 득점기록자의 판정들은 최종적이고 무오류적(infallible)일
것이다. 아니 오히려 그것들이 오류 가능적인지 무오류적인지를 묻는 질문은
무의미할 것이다. 그에게 ‘옳게’ 또는 ‘그르게’ 할 것이 아무것도 없기
때문이다. 그러나 보통의 게임에서 ‘득점은 득점기록자가 말하는 것이다’는
득점 규칙이 아니다. 그것은 특정 사례들에서 득점 규칙에 대한 그의 적용의
권위와 최종성을 마련하는 규칙이다.

권위 있는 결정(authoritative decision)의 이 예에서 배울 두 번째 교훈은
더 근본적인 문제들에 닿아 있다. 우리는 보통의 게임과 ‘득점기록자의 재량’
게임을 구별할 수 있는데, 이는 득점 규칙이 다른 규칙들과 마찬가지로
득점기록자가 선택을 행사해야 하는 개방적 직조(open texture)의 영역을
가지고 있음에도 불구하고, 확정된 의미의 핵심(core of settled meaning)을
가지고 있기 때문이다. 바로 이것이 득점기록자가 자유롭게 이탈할 수 없는
것이며, 그것이 미치는 한에서 선수의 득점에 관한 비공식 진술들과
득점기록자의 공식 판정들 모두에 대해 올바른 득점과 올바르지 않은 득점의
기준을 구성한다. 바로 이것이 득점기록자의 판정들은 비록 최종적이지만
무오류적이지는 않다고 말하는 것을 참으로 만든다. 법에서도 동일하다.

어느 지점까지는 득점기록자가 내린 몇몇 판정들이 명백히 그르다는 사실은
게임이 계속되는 것과 양립하지 않는 것이 아니다. 그것들은 명백히 옳은
판정들만큼이나 동등하게 산입된다(count as much). 그러나 올바르지 않은
결정들(incorrect decisions)에 대한 관용이 동일한 게임의 계속된 존재와
양립할 수 있는 정도에는 한계가 있으며, 이것은 중요한 법적 유비를 가진다.
고립적이거나 예외적인 공식적 일탈들(official aberrations)이 용인된다는
사실은 크리켓이나 야구가 더 이상 행해지고 있지 않다는 것을 의미하지
않는다. 다른 한편, 이러한 일탈들이 빈번하거나 득점기록자가 득점 규칙을
거부한다면, 어느 지점에는 반드시 선수들이 득점기록자의 일탈적 판정들을
더 이상 수용하지 않거나, 만약 그들이 수용한다면 게임이 달라진 지점이
오게 된다. 그것은 더 이상 크리켓이나 야구가 아니라 ‘득점기록자의
재량’이다. 왜냐하면 이러한 다른 게임들의 규정적 특징(defining
feature)은, 그 개방적 직조가 득점기록자에게 어떤 \textbf{(p.~145)}
재량의 폭(latitude)을 남겨두든, 일반적으로 그 결과들이 규칙의 명백한
의미(plain meaning)가 요구하는 방식으로 평가되어야 한다는 점이기
때문이다. 어떤 상상 가능한 조건에서는 실제로 행해지고 있는 게임이
‘득점기록자의 재량’이라고 말해야 할 것이다. 그러나 모든 게임들에서
득점기록자의 판정들이 최종적이라는 사실은 모든 게임들이 그것이라는 것을
의미하지 않는다.

특정 사례에서 법이 무엇인지에 대한 최종적이고 권위 있는 진술로서 법원의
결정이 가지는 고유한 지위에 근거한 규칙회의주의(rule-scepticism)의
형태를 평가할 때에는 이러한 구별들을 마음에 두어야 한다. 법의 개방적
직조는 법원들에게, 그 결정들이 법창출 선례들(law-making precedents)로
사용되지 않는 득점기록자들에게 남겨진 것보다 훨씬 더 넓고 중요한 법창출
권한(law-creating power)을 남긴다. 법원들이 무엇을 결정하든, 그것이
모두에게 명백해 보이는 규칙 부분 안에 놓인 문제이든 논쟁 가능한 경계에
놓인 문제이든, 입법에 의해 변경될 때까지 존립한다. 그리고 그 입법의
해석에 대해서도 법원들은 다시 동일한 최후의 권위 있는 목소리를 가질
것이다. 그럼에도 불구하고, 법원 체계를 세운 뒤 법은 최고 법원이
적합하다고 생각하는 것이면 무엇이든 된다고 규정하는 헌법과, 실제의 미국
헌법---또는 그 문제에 관해서는 어떤 현대 국가의 헌법---사이에는 여전히
구별이 남아 있다. ‘헌법(또는 법)은 판사들이 그것이라고 말하는 것이
무엇이든 그것이다’가 이 구별을 부정하는 것으로 해석된다면, 그것은
거짓이다. 어느 특정 시점에서 판사들, 심지어 최고 법원의 판사들도 그
규칙들이 중심에서는 올바른 사법적 결정의 기준들을 제공할 만큼 충분히
확정적인 체계의 일부이다. 이러한 기준들은 법원들에 의해, 체계 안에서
다툴 수 없는 결정들을 내릴 권위를 행사할 때 자유롭게 무시할 수 없는
것으로 간주된다. 어떤 개별 판사도 자기 직무에 올 때, 어떤 득점기록자가
자기 직무에 올 때와 마찬가지로, 의회 내 여왕의 제정물들(enactments)은
법이라는 규칙과 같은 규칙이 전통으로 확립되어 있고 그 직무 수행의
기준으로 수용되어 있음을 발견한다. 이것은 그 직무 점유자들의 창조적
활동을 허용하면서도 경계 짓는다. 이러한 기준들은 기실 그 시대 판사들
대부분이 그것들을 고수하지 않는다면 계속 존재할 수 없는데, 특정 시점에서
그것들의 존재는 단순히 올바른 재판의 기준들로 그것들이 수용되고 사용되는
데 있기 때문이다. 그러나 이것이 그 기준들을 사용하는 판사를 그 기준들의
창안자(author)로 만들거나, 호들리(Hoadly)의 언어로 말해 자기가 원하는
대로 \textbf{(p.~146)} 결정할 권능을 가진 ‘입법자(lawgiver)’로 만들지는
않는다. 기준들을 유지하기 위해 판사의 고수는 필요하지만, 판사가 그것들을
만드는 것은 아니다.

물론 사법적 결정들을 최종적이고 권위 있는 것으로 만드는 규칙들의
보호막(shield) 배후에서, 판사들이 결합하여 기존 규칙들을 거부하고 심지어
가장 명확한 의회 법률들(Acts of Parliament)조차도 자신들의 결정들에 어떤
한계도 부과하는 것으로 간주하기를 중단할 가능성은 있다. 그들의 판정들
대다수가 이러한 성격을 가지고 수용된다면, 이것은 크리켓에서
‘득점기록자의 재량’으로 게임이 전환되는 것과 평행한 체계의 변형에 해당할
것이다. 그러나 그러한 변형들의 상존 가능성은, 그 변형이 일어났다면
체계가 그러했을 바로 그것이 현재의 체계라는 것을 보여 주지 않는다. 어떤
규칙들도 위반이나 거부(repudiation)에 대해 보장될 수 없다. 인간 존재들이
그것들을 깨거나 거부하는 것은 결코 심리적으로나 물리적으로 불가능하지
않기 때문이다. 그리고 충분히 많은 사람들이 충분히 오래 그렇게 한다면, 그
규칙들은 존재를 중단할 것이다. 그러나 특정 시점에서 규칙들이 존재한다는
것은 파괴에 대한 이러한 불가능한 보장들이 있어야 한다는 것을 요구하지
않는다. 어떤 시점에 판사들이 의회 법률들(Acts of Parliament) 또는
연방의회 법률들(Acts of Congress)을 법으로 수용할 것을 요구하는 규칙이
있다고 말하는 것은, 첫째, 이 요구사항에 대한 일반적 준수(compliance)가
있고 개별 판사들의 일탈이나 거부는 드물다는 것, 둘째, 그것이 발생하거나
발생한다면 우세한 다수(preponderant majority)에 의해 심각한 비판의
대상으로, 그리고 그른(wrong) 것으로 취급되거나 취급될 것이라는 것을
수반한다. 비록 특정 사례에서 그에 따라 내려진 결정의 결과는, 결정들의
최종성에 관한 규칙 때문에, 그 결정의 올바름(correctness)이 아니라
유효성(validity)을 인정하는 입법에 의해서만 상쇄될 수 있다고 하더라도
말이다. 인간 존재들이 자기 약속들을 모두 어기는 것은 논리적으로
가능하다. 처음에는 아마도 그것이 그른 행위라는 감각을 가지고, 그다음에는
그러한 감각 없이 그럴 수 있다. 그러면 약속을 지키는 것을 의무적인 것으로
만드는 규칙은 존재를 중단할 것이다. 그러나 이것은 현재 그러한 규칙이
존재하지 않고 약속들이 실제로 구속력이 없다는 견해를 빈약하게 지지할
뿐이다. 판사들의 경우에서, 그들이 현재 체계의 파괴를 기획해 이끌어낼
가능성에 근거한 평행 논변은 이보다 더 강하지 않다.

규칙회의주의의 주제를 떠나기 전에 우리는, 규칙들이 법원들의 결정들에
대한 예측들이라는 그 적극적 논지(contention)에 관해 마지막 말을 해야
한다. 여기에 어떤 진리가 있든, 그것은 기껏해야 사적 개인들이나 그
조언자들이 제시하는 법 진술들(statements of law)에 적용될 수 있다는 점은
명백하며 자주 \textbf{(p.~147)} 언급되어 왔다. 그것은 법원들 자신의
법규칙 진술들에는 적용될 수 없다. 이것들은 일부 더 극단적인
‘현실주의자들(Realists)’이 주장했듯이 속박 없는 재량 행사를 위한 언어적
덮개(verbal covering)이거나, 아니면 법원들이 내부 관점에서 올바른 결정의
기준으로 진정으로 간주하는 규칙들의 정식화여야 한다. 다른 한편, 사법적
결정들에 대한 예측들은 법 안에서 부정할 수 없이 중요한 자리를 가진다.
개방적 직조(open texture)의 영역에 도달할 때, ‘이 문제에서 법은
무엇인가?’라는 질문에 대한 답으로 우리가 유익하게 제공할 수 있는 전부가
법원들이 무엇을 할지에 대한 조심스러운 예측인 경우가 매우 흔하다. 더
나아가, 규칙들이 요구하는 것이 모두에게 명확한 경우에도 그것에 대한
진술은 법원의 결정에 대한 예측의 형식으로 이루어지는 경우가 많다. 그러나
주로 후자의 경우에, 그리고 전자의 경우에는 다양한 정도로, 그러한 예측의
토대는 법원들이 법규칙들을 예측들로서가 아니라 결정에서 따라야 할
기준들로 간주한다는 지식이라는 점을 주목하는 것이 중요하다. 그 기준들은
그 개방적 직조에도 불구하고 법원의 재량을 배제하지는 않더라도 제한할
만큼 충분히 확정적이다. 따라서 많은 경우 법원이 무엇을 할지에 대한
예측들은, 체스 선수들이 비숍을 대각선으로 움직일 것이라고 우리가 할 수
있는 예측과 같다. 그것들은 궁극적으로 규칙들의 비예측적(non-predictive)
측면과, 그 예측들이 관련되는 자들이 기준들로서 수용한 규칙들의 내부
관점에 대한 이해에 기초한다(rest on). 이것은, 어떤 사회 집단에서
규칙들이 존재한다는 사실이 예측들을 가능하게 하고 자주 신뢰할 수 있게
만들지만 그것과 동일시될 수는 없다는, 제5장에서 이미 강조한 사실의 또
다른 측면일 뿐이다.

\subsection{4. 승인 규칙에서의
불확실성(UNCERTAINTY)}\label{uxc2b9uxc778-uxaddcuxce59uxc5d0uxc11cuxc758-uxbd88uxd655uxc2e4uxc131uncertainty}

형식주의와 규칙회의주의는 법리학적 이론(juristic theory)의
스킬라(Scylla)와 카립디스(Charybdis)이다. 그것들은 서로를 교정하는
곳에서는 유익한 거대한 과장들이며, 진리는 그 사이에 있다. 이 중도적 길을
유익한 세부 내용으로 특징짓고, 제정법(statute)이나 선례에서 법의 개방적
직조가 법원들에게 남겨 두는 창조적 기능을 수행할 때 법원들이 특징적으로
사용하는 다양한 추론 유형들을 보여 주기 위해서는, 기실 여기에서 시도될
수 없는 많은 일이 행해져야 한다. 그러나 우리는 \textbf{(p.~148)} 이
장에서, 제6장 말미에 미뤄 두었던 중요한 주제로 유익하게 돌아갈 수 있을
만큼은 말했다. 이것은 특정 법규칙들의 불확실성이 아니라 승인 규칙의
불확실성, 따라서 법원들이 유효한 법규칙들을 식별할 때 사용하는 궁극적
기준들의 불확실성에 관한 것이었다. 특정 규칙의 불확실성과 그것을 체계의
규칙으로 식별하는 데 사용되는 기준의 불확실성 사이의 구별은 그 자체로
모든 경우에 명확한 것은 아니다. 그러나 규칙들이 권위 있는 텍스트를 가진
제정법 형식의 제정물들(statutory enactments)인 경우에 그 구별은 가장
명확하다. 제정법(statute)의 문언과 그것이 특정 사례에서 요구하는 것은
완전히 명백할 수 있다. 그런데도 입법부가 이런 방식으로 입법할
권한(power)을 가지는지에 관해서는 의문들이 있을 수 있다. 때때로 이러한
의문들의 해소(resolution)는 입법 권한을 부여한 또 다른 법규칙의 해석만을
요구하며, 이 규칙의 유효성은 의심스럽지 않을 수 있다. 예컨대 종속된
권한기관(subordinate authority)이 만든 제정물(enactment)의 유효성이
문제되는 경우가 그럴 것이다. 종속된 권한기관의 입법 권한들을 규정하는
모법(parent Act of Parliament)의 의미에 관해 의문들이 발생하기 때문이다.
이것은 특정 제정법(statute)의 불확실성 또는 개방적 직조의 사례일 뿐이며
어떤 근본적 문제도 제기하지 않는다.

그러한 통상적 문제들과 구별되어야 할 것은 최고 입법부(supreme
legislature) 자체의 법적 권능(legal competence)에 관한 문제들이다.
이것들은 법적 유효성의 궁극적 기준들(ultimate criteria of legal
validity)에 관한 것이며, 최고 입법부의 권능을 특정하는 성문헌법(written
constitution)이 없는 우리 자신의 체계와 같은 법체계에서도 발생할 수
있다. 압도적 다수의 사례들에서 ‘의회 내 여왕이 제정하는 것은 무엇이든
법이다’라는 공식은 의회의 법적 권능에 관한 규칙의 적절한 표현이며,
그렇게 식별된 규칙들이 그 주변부에서는 아무리 열려 있을 수 있더라도,
법의 식별을 위한 궁극적 기준으로 수용된다. 그러나 그 의미나 범위에 관해
의문들이 발생할 수 있다. 우리는 ‘의회에 의해 제정된(enacted by
Parliament)’이라는 말이 무엇을 의미하는지 물을 수 있으며, 의문들이
발생하면 그것들은 법원들에 의해 해소될 수 있다. 한 법체계의 궁극적
규칙이 이처럼 의심스러울 수 있고 법원들이 그 의문을 해소할 수 있다는
사실로부터, 법체계 안에서 법원들의 위치에 관해 어떤 추론이 도출되어야
하는가? 그것은 법체계의 토대가 법적 유효성의 기준들을 특정하는 수용된
승인 규칙이라는 테제에 어떤 한정을 요구하는가?

\textbf{(p.~149)} 이러한 질문들에 답하기 위해 여기서는 영국
의회주권(parliamentary sovereignty) 교설의 몇몇 측면들을 검토할 것이다.
물론 유사한 의문들은 어떤 체계에서든 법적 유효성의 궁극적 기준들과
관련하여 발생할 수 있다. 법은 본질적으로 법적으로 제약받지 않는 의지의
산물이라는 오스틴주의 교설의 영향 아래에서, 오래된 헌법이론가들은 마치
주권적인 입법부가 존재해야 한다는 것이 논리적 필연성인 것처럼 썼다.
여기서 주권적이라는 것은, 그 입법부가 계속적 기구(continuing body)로
존재하는 모든 순간에 \emphksans{외부로부터(ab extra)} 부과된 법적 제한들뿐
아니라 자기 자신의 선행 입법(prior legislation)으로부터도 자유롭다는
의미이다. 이러한 의미에서 의회가 주권적이라는 점은 이제 확립된 것으로
간주될 수 있으며, 이전 의회는 자기 입법을 폐지하지 못하도록 자신의
‘승계자들(successors)’을 막을 수 없다는 원칙은 법원들이 유효한
법규칙들을 식별할 때 사용하는 궁극적 승인 규칙의 일부를 구성한다. 그러나
그러한 의회가 존재해야 한다고 명하는 것은 논리의 필연성도, 하물며 자연의
필연성도 아니라는 점을 보는 것이 중요하다. 그것은 똑같이 상상 가능한
여러 배치들 가운데 하나일 뿐이며, 우리의 체계에서 법적 유효성의 기준으로
수용되게 된 것이다. 그러한 다른 배치들 가운데에는 아마도 더 나을지도
모르게, 똑같이 ‘주권’이라는 이름을 받을 자격이 있는 또 다른 원칙이 있다.
그것은 의회가 자신의 승계자들의 입법 권능(legislative competence)을
불가역적으로 제한할 수 없는 상태여서는 \emphksans{안 되며(not)}, 반대로
이러한 더 넓은 자기제한 권한(self-limiting power)을 가져야 한다는
원칙이다. 그러면 의회는 그 역사에서 적어도 한 번은 수용되고 확립된
교설이 허용하는 것보다 훨씬 더 큰 입법 권능의 영역을 행사할 능력을
가지게 될 것이다. 의회가 그 존재의 모든 순간에, 심지어 자기 자신에 의해
부과된 것들을 포함한 법적 제한들로부터 자유로워야 한다는 요구는, 결국
법적 전능(legal omnipotence)이라는 애매한 관념의 한 해석일 뿐이다.
그것은 사실상, 잇따른 의회들(successive parliaments)의 입법 권능에
영향을 미치지 않는 모든 사항에서의 \emphksans{계속적(continuing)} 전능과, 그
행사가 오직 한 번만 향유될 수 있는 제한 없는
\emphksans{자기포섭적(self-embracing)} 전능 사이에서 선택한다. 이 두 전능
구상들은 전능한 신에 관한 두 구상에서 평행을 가진다. 한편에는 자기
존재의 모든 순간에 동일한 능력들(powers)을 누리고 따라서 그 능력들을
줄일 수 없는 신이 있고, 다른 한편에는 장래를 향해 자기 전능을 파괴할
능력(power)을 포함하는 능력들을 가진 신이 있다. 우리의 의회가 어떤
형태의 전능---계속적 전능인지 자기포섭적 전능인지---을 누리는지는 법을
식별하는 궁극적 기준으로 수용되는 규칙의 형식에 관한 경험적 질문이다.
그것은 법체계의 토대에 놓인 규칙에 관한 질문이지만, 여전히 어느 주어진
시점에는 적어도 어떤 점들에 관해 상당히 확정적인 답이 있을 수 있는
사실의 문제이다. 따라서 현재 수용된 규칙은 계속적 주권의 규칙이며,
그래서 의회는 자기 제정법들(statutes)을 폐지로부터 보호할 수 없다는 점은
명확하다.

그러나 다른 모든 규칙에서와 마찬가지로, 의회주권(parliamentary
sovereignty)의 규칙이 이 지점에서 확정적(determinate)이라는 사실은
그것이 모든 지점에서 그러하다는 것을 의미하지 않는다. 현재 명확히 옳거나
그른 답이 없는 질문들이 그것에 관해 제기될 수 있다. 이것들은 이 문제에서
결국 그 선택들에 권위가 부여되는 누군가가 하는 선택에 의해서만 확정될 수
있다. 의회주권 규칙 안의 그러한 불확정성들(indeterminacies)은 다음
방식으로 나타난다. 현행 규칙 아래에서는, 의회가 제정법(statute)에 의해
어떤 사항을 장래 의회의 입법 범위에서 불가역적으로 철회할 수 없다는 점이
인정된다. 그러나 단순히 그렇게 한다고 표방하는 제정물(enactment)과,
의회가 여전히 어떤 사항에 관해서든 입법할 가능성을 열어 두면서도 입법의
‘방식과 형식(manner and form)’을 변경한다고 표방하는 제정물 사이에는
구별이 그어질 수 있다. 후자는 예컨대 특정 쟁점들에 대해서는 양원이 함께
앉아 다수결로 통과시키거나 국민투표(plebiscite)로 확인되지 않는 한 어떤
입법도 법으로서 효력을 갖지 못한다고 요구할 수 있다. 그것은 그 조항
자체가 동일한 특별 절차에 의해서만 폐지될 수 있다는 규정으로 그러한
조항을 ‘강화 고착(entrench)’할 수 있다. 입법 과정에서의 그러한 부분적
변경은 의회가 자기 승계자들을 불가역적으로 구속할 수 없다는 현행 규칙과
충분히 양립할 수 있다. 그것이 하는 일은 승계자들을
\emphksans{구속(bind)}한다기보다는, 특정 쟁점들에 \emphksans{관한 한(quoad)}
그들을 제거하고 이 쟁점들에 관한 그들의 입법 권한들을 새로운 특별
기구(body)로 이전하는 것이기 때문이다. 따라서 이러한 특별 쟁점들과
관련하여 의회는 의회를 ‘구속’하거나 ‘족쇄를 채우거나(fettered)’ 그
계속적 전능을 감소시킨 것이 아니라, ‘의회’와 입법하기 위해 행해져야 하는
것을 ‘재규정(redefined)’했다고 말할 수 있다.

명백히 이 장치가 유효하다면, 의회는 \textbf{(p.~151)} 그것을 사용하여
의회가 자기 승계자들을 구속할 수 없다는 수용된 교설이 자신의 권한 밖에
두는 것처럼 보이는 결과들과 매우 유사한 결과들을 달성할 수 있을 것이다.
왜냐하면 의회가 입법할 수 있는 영역을 경계 짓는 것과 입법의 방식과
형식만을 변경하는 것 사이의 차이는 기실 몇몇 경우에는 충분히 명확하지만,
실제 효과상 이 범주들은 서로 점차 맞물리며 경계가 흐려지기
때문이다(shade into each other). 엔지니어들의 최저임금을 정한 뒤,
엔지니어 임금에 관한 어떤 법안도 엔지니어 노동조합(Engineers' Union)의
결의로 확인되지 않는 한 법으로서 효력을 가지지 않는다고 규정하고, 이어서
이 조항을 강화 고착한 제정법(statute)은, 임금을 ‘영원히’ 고정하고 그
폐지를 거칠게 전면 금지한 제정법(statute)이 실제로 할 수 있는 모든 것을
기실 확보할 수도 있다. 그러나 법률가들이 어느 정도 설득력을 가진 것으로
인정할 논변(argument)은, 비록 후자는 현재의 계속적 의회주권 규칙
아래에서는 효력이 없을 것이지만 전자는 그렇지 않을 것임을 보이도록
제시될 수 있다. 이 논변의 단계들은 의회가 무엇을 할 수 있는지에 관한
일련의 논지들(contentions)로 이루어지며, 각 논지는 그 앞의 논지와 일정한
유비를 가지면서도 그 앞의 논지보다 더 적은 동의를 얻을 것이다. 그중 어떤
것도 그르다고 배제될 수 없고, 옳다고 자신 있게 수용될 수도 없다. 우리는
체계의 가장 근본적인 규칙의 개방적 직조 영역 안에 있기 때문이다.
여기서는 어느 순간에도 답은 없고 오직 답들만 있는 질문이 제기될 수 있다.

따라서 의회가 상원을 완전히 폐지함으로써 의회의 현재 구성을 불가역적으로
변경할 수도 있고, 그래서 특정 입법에 대한 상원의 동의를 면제했으며 일부
권위 있는 논자들이 단지 여왕과 하원에 대한 의회 권한들 일부의 철회
가능한 위임으로 해석하기를 선호하는 1911년 및 1949년 의회법들(Parliament
Acts)을 넘어설 수 있다는 점이 인정될지도 모른다. 또한 다이시(Dicey)가
단언했듯이, \footnote{\emph{The Law of the Constitution}(10th edn.),
  p.~68 n.} 의회가 자기 권한들이 끝났음을 선언하고 장래 의회들의 선거를
마련하는 입법을 폐지하는 법률(Act)에 의해 자기 자신을 완전히 파괴할 수
있다는 점도 인정될지도 모른다. 그렇다면 의회는 맨체스터 법인(Manchester
Corporation) 같은 어떤 다른 기관에 자신의 모든 권한들을 이전하는 법률을
통해 이러한 입법적 자살(legislative suicide)을 유효하게 수반시킬 수도
있을 것이다. 의회가 이것을 할 수 있다면, 그보다 적은 것을 실효적으로 할
수는 없는가? 특정 사항들에 관해 \textbf{(p.~152)} 입법할 자신의 권한들을
끝내고, 그것들을 자기 자신과 어떤 추가 기관을 포함하는 새로운 복합
실체(composite entity)에 이전할 수는 없는가? 이러한 기초 위에서,
도미니언(Dominion)에 영향을 미치는 어떤 입법에 대해서도 도미니언의
동의를 요구한 웨스트민스터 제정법(Statute of Westminster) 제4조가,
도미니언을 위해 입법할 의회의 권한들과 관련하여 실제로 이 일을 했을 수도
있지 않은가? 이 조항이 도미니언의 동의 없이 실효적으로 폐지될 수 있다는
논지(contention)는, 상키 경(Lord Sankey)이 말했듯이 ‘현실과 아무 관계가
없는’ ‘이론’일 뿐만 아니라, 나쁜 이론일 수도 있다. 또는 적어도 반대
이론보다 나을 것이 없을 수 있다. 마지막으로, 의회가 자기 자신의 행위에
의해 이런 방식들로 재구성될 수 있다면, 왜 그것은 엔지니어
노동조합(Engineers' Union)이 특정 유형의 입법에서 필수적인 동의
요소(consenting element)가 되도록 규정함으로써 자기 자신을 재구성할 수
없는가?

이 논변에서 의심스럽지만 명백히 잘못된 것은 아닌 단계들을 구성하는 몇몇
의문스러운 명제들이, 언젠가 그 문제를 결정하도록 요청받은 법원에 의해
승인되거나 거부될 가능성은 충분히 있다. 그러면 우리는 그것들이 제기하는
질문들에 대한 답을 갖게 될 것이며, 그 답은 체계가 존재하는 한, 주어질 수
있는 답들 가운데 유일한 권위 있는 지위를 가질 것이다. 법원들은 이
지점에서 유효한 법을 식별하는 궁극적 규칙을 확정적인 것으로 만들었을
것이다. 여기서 ‘헌법은 판사들이 말하는 것이다’는 말은 최고
법원들(supreme tribunals)의 특정 결정들이 다투어질 수 없다는 것만을
의미하지 \emphksans{않는다}. 처음 보기에는 그 광경이 역설적으로 보인다.
여기에는 법원들이, 판사로서 자신들에게 관할권(jurisdiction)을 부여하는
바로 그 법들의 유효성 자체가 그 기준들에 의해 검사되어야 하는 궁극적
기준들을 확정하는 창조적 권한들을 행사하고 있기 때문이다. 어떻게 헌법이
헌법이 무엇인지를 말할 권위를 부여할 수 있는가? 그러나 모든 규칙이 어떤
지점들에서는 의심스러울 수 있지만, 모든 규칙이 모든 지점에서 의심에 열려
있지는 않다는 것이 기실 법체계가 존재하기 위한 필요조건이라는 점을
기억하면 이 역설은 사라진다. 어느 주어진 시점에 법원들이 궁극적 유효성
기준들에 관한 이러한 한계 질문들(limiting questions)을 결정할 권위를
가질 가능성은, 단지 그 시점에는 그 권위를 부여하는 규칙들을 포함하여
광대한 법 영역에 대한 그 기준들의 적용이 아무 의문도 제기하지 않는다는
사실에 의존한다. 비록 그 기준들의 정확한 범위(scope)와 적용범위(ambit)에
대해서는 의문이 제기되더라도 말이다.

\textbf{(p.~153)} 그러나 이 답은 어떤 이들에게는 그 문제를 너무 쉽게
처리하는 것처럼 보일 수 있다. 그것은 법적 유효성의 기준들을 특정하는
근본 규칙들(fundamental rules)의 가장자리에서 법원들이 하는 활동을 매우
부적절하게 특징짓는 것처럼 보일 수 있다. 그 이유는 아마도 그것이 이
활동을, 불확정적인 것으로 드러난 특정 제정법(statute)을 해석하면서
법원들이 창조적 선택을 행사하는 통상적 사례들에 지나치게 가깝게
동화시키기 때문일 수 있다. 그러한 통상적 사례들은 어떤 체계에서든
발생해야 한다는 점은 명확하며, 그래서 법원들이 그 사례들을, 제정법이
열어 둔 대안들 사이에서 선택함으로써 확정할 관할권을 가진다는 점은,
그들이 이 선택을 발견(discovery)으로 위장하기를 선호하더라도, 법원들이
근거하여 행위하는 규칙들의 일부, 비록 암묵적 일부일 뿐이더라도, 명백히
일부인 것처럼 보인다. 그러나 적어도 성문헌법(written constitution)이
없는 경우에는, 유효성의 근본 기준들에 관한 질문들은 종종 이렇게 사전에
염두에 둘 수 있는 성질을 가지지 \emphksans{않는(not)} 것처럼 보인다. 그리고
바로 이러한 사전에 염두에 둘 수 있는 성질이 있어야, 그러한 질문들이
발생할 때 법원들이 이미 기존 규칙들 아래에서 이 종류의 질문들을 확정할
명확한 권위를 가지고 있다고 말하는 것이 자연스럽다.

‘형식주의적(formalist)’ 오류의 한 형태는 어쩌면 법원이 취하는 모든
단계가 그것을 취할 권위를 사전에 부여하는 어떤 일반 규칙에 의해 포괄되어
있어서, 법원의 창조적 권한들이 \emphksans{언제나(always)} 위임된 입법
권한(delegated legislative power)의 한 형태라고 생각하는 바로 그것일 수
있다. 진실은, 법원들이 가장 근본적인 헌법 규칙들에 관한 사전에 염두에
두어지지 않았던 질문들을 확정할 때, 그 질문들이 발생하고 결정이 내려진
뒤에야 그 질문들을 결정할 권위가 수용되게 된다는 것일 수 있다. 여기서는
성공하는 것은 오직 성공뿐이다(all that succeeds is success). 문제된
헌법적 질문이 사회를 너무 근본적으로 분열시켜 사법적 결정에 의한 처리를
허용하지 않을 수도 있다는 것은 상상 가능하다. 1909년 남아프리카법(South
Africa Act)의 고착 조항들(entrenched clauses)에 관한 남아프리카의
쟁점들은 한때 법적 처리를 하기에는 너무 분열적일 위험이 있었다. 그러나
덜 중요한 사회적 쟁점들이 관련된 곳에서는, 법의 바로 그 원천들에 관한
매우 놀라운 사법적 법창출(judicial law-making)도 조용히
‘삼켜질(swallowed)’ 수 있다. 그러한 경우에는,
\emphksans{회고적으로(retrospect)} 법원들이 자신들이 한 일을 할 ‘내재적’
권한(inherent power)을 언제나 가지고 있었다고 자주 말해질 것이고,
진정으로 그렇게 보일 수도 있다. 그러나 그 유일한 증거가 행해진 일의
성공뿐이라면, 이것은 경건한 허구(pious fiction)일 수 있다.

영국 법원들이 선례의 구속력(binding force)에 관한 규칙들을 조작한 것은,
\textbf{(p.~154)} 아마도 이 마지막 방식으로, 권한들을 확보하여
행사하려는 성공적인 시도(successful bid)로서 서술되는 것이 가장 정직할
것이다. 여기서 권한(power)은 성공으로부터 \emphksans{사후적으로(ex post
facto)} 권위(authority)를 획득한다. 따라서 형사항소법원(Court of
Criminal Appeal)이 \emph{Rex} 대 \emph{Taylor} 사건 \footnote{{[}1950{]}
  2 KB 368.}에서 신민의 자유에 관한 사항들에 대해 자기 자신의 선례들에
구속되지 않는다고 판정할 권위를 가진다고 판시하기 전에는, 그 법원이
그러한 권위를 가지는지는 전적으로 열려 있는 것으로 보였을 수 있다.
그러나 그 판정은 내려졌고, 지금은 법으로서 따라지고 있다. 법원이 이런
방식으로 판정할 내재적 권한을 언제나 가지고 있었다는 진술은, 분명 상황을
실제보다 더 정돈되어 보이게 하는 방식일 뿐일 것이다. 여기, 바로 이러한
근본적인 것들의 가장자리에서는, 우리는 규칙회의주의자를 환영해야 한다.
다만 그는 자신이 환영받는 곳이 가장자리라는 점을 잊지 않아야 하며,
법원들이 가장 근본적인 규칙들을 이렇게 눈에 띄게 발전시킬 수 있게 하는
것은 상당 부분 법의 광대한 중심 영역들에서 의심할 여지 없이
규칙-지배적인(rule-governed) 작용들을 해 온 데서 법원들이 모은
위신(prestige)이라는 사실을 우리에게 보지 못하게 해서는 안 된다.

\newpage

\subsection{제7장 주석}\label{uxc81c7uxc7a5-uxc8fcuxc11d}


{\footnotesize
\noteentry 125쪽. \notetitle{규칙의 예시적 전달(Communication of rules by examples).} 이러한 용어로 선례(precedent)의 사용을 특징지은 설명으로는 Levi의 “An
Introduction to Legal Reasoning” 제1절, \emph{University of Chicago Law
Review} 제15권(1948년)을 보라. Wittgenstein은 \emph{Philosophical
Investigations} (특히 제1부 제208--238절)에서 규칙을 가르치고 따르는
일에 관한 관념들(notions)에 대해 여러 중요한 통찰을 제시한다.
Wittgenstein에 대한 논의는 Winch, \emph{The Idea of a Social Science}
제24--33쪽, 91--93쪽 참조.

\noteentry 128쪽. \notetitle{언어적으로 정식화된 규칙들의 개방적 직조(Open texture of verbally formulated rules).} 개방적 직조(open texture)라는 관념에
대해서는 Waismann의 “Verifiability”, \emph{Essays on Logic and
Language} 제1권(Flew 편집), 제117--130쪽을 보라. 법적 추론과 관련한
논의로는 Dewey, “Logical Method and Law,” \emph{Cornell Law Quarterly}
제10권(1924); Stone, \emph{The Province and Function of Law}, 제6장;
Hart, “Theory and Definition in Jurisprudence,” \emph{Proceedings of
the Aristotelian Society} 보충권 제29권(1955), 제258--264쪽;
“Positivism and the Separation of Law and Morals,” \emph{Harvard Law
Review} 제71권(1958), 제606--612쪽 참조.

\noteentry 129쪽. \notetitle{형식주의(formalism)와 개념주의(conceptualism).} 법
문헌들에서 이 용어들과 거의 유의어로 사용되는 표현으로는 ‘기계적 또는
자동적 법리학(mechanical or automatic jurisprudence)’, ‘구상들의
법리학(jurisprudence of conceptions)’, ‘논리의 과도한 사용(excessive use
of logic)’ 등이 있다. 이와 관련하여 Pound, “Mechanical Jurisprudence,”
\emph{Columbia Law Review} 제8권(1908); \emph{Interpretations of Legal
History}, 제6장 참조. 이러한 용어들로 정확히 어떤 결함 또는 폐해가
지칭되는지는 항상 명확하지 않다. 이에 관해서는 Jensen, \emph{The Nature
of Legal Argument}, 제1장, 그리고 Honoré의 서평, \emph{Law Quarterly
Review} 제74권(1958), 제296쪽; Hart, 위에 언급한 \emph{Harvard Law
Review} 제71권, 제608--612쪽을 참조.

\noteentry 131쪽. \notetitle{법적 기준들과 구체적 규칙들(Legal standards and specific rules).} 이러한 두 가지 법적 통제 형태의 성격과 상호 관계에 대한 가장
통찰력 있는 일반적 논의는 Dickinson, \emph{Administrative Justice and
the Supremacy of Law}, 제128--140쪽에 있다.

\noteentry 131쪽. \notetitle{행정 규칙제정에 의해 구체화되는 법적 기준들(Legal standards implemented by administrative rule-making).} 미국에서는 Interstate
Commerce Commission, Federal Trade Commission 등 연방 규제기관들이 ‘공정
경쟁(fair competition)’, ‘공정하고 합리적인 요율(just and reasonable
rates)’ 등의 광범위한 법적 기준들을 구체화하는 규칙들을 제정한다.
(Schwartz, \emph{An Introduction to American Administrative Law},
제6--18쪽, 33--37쪽 참조.) 영국에서도 이와 유사한 규칙제정 기능이
행정부에 의해 수행되지만, 일반적으로 미국에서 익숙한 이해관계자들에 대한
형식적 준사법적 청문(quasi-judicial hearing) 절차 없이 이루어진다.
예컨대, 1957년 \emph{Factories Act} 제46조에 근거한 복지규정(Welfare
Regulations), 같은 법 제60조에 의한 건축규정(Building Regulations)을
참조. 또한, 1947년 \emph{Transport Act}에 따라 이의제기자들(objectors)의
의견을 청취한 후 ‘요금 체계(charges scheme)’를 확정할 수 있는
\emph{Transport Tribunal}의 권한은 미국식 모델에 보다 근접해 있다.

\noteentry 132쪽. \notetitle{주의 기준들(Standards of care).} 주의의무(duty of care)의
구성요소들에 대한 통찰력 있는 분석으로는 Learned Hand 판사의 판결문
\emph{US v. Carroll Towing Co.} (1947), 159 F 2nd 169, 173쪽을 보라.
일반적 기준들을 구체적 규칙들로 대체하는 바람직성에 대해서는 Holmes,
\emph{The Common Law}, 제3강, 제111--119쪽 참조. 이 견해에 대한 비판은
Dickinson, 앞의 책, 제146--150쪽에 제시되어 있다.

\noteentry 133쪽. \notetitle{구체적 규칙들에 의한 통제(Control by specific rules).} 유연한 기준들(flexible standards)보다는 엄격하고 고정된 규칙들(hard and
fast rules)이 적절한 통제 형태가 되는 조건에 대해서는 Dickinson, 앞의
책, 제128--132쪽, 제145--150쪽을 보라.

\noteentry 134쪽. \notetitle{선례와 법원의 입법 활동(Precedent and the legislative activity of Courts).} 영국법에서의 선례(precedent in English law)의
현대적 일반 설명으로는 R. Cross, \emph{Precedent in English Law}
(1961)을 보라. 본문에서 언급된 ‘협소화 과정(narrowing process)’의 잘
알려진 예시는 \emph{L. \& S. W. Railway Co.} v. \emph{Gomm} (1880), 20
Ch.D. 562로, 이는 \emph{Tulk} v. \emph{Moxhay} (1848), 2 Ph. 774에서의
법리를 좁혀 적용하였다.

\noteentry 136쪽. \notetitle{규칙회의주의의 다양상들(Varieties of rule-scepticism).} 이
주제에 대한 미국 법학의 논의는 하나의 논쟁(debate)으로 읽을 때 특히
유익하다. 예컨대 \emph{Law and the Modern Mind} (특히 제1장 및 부록 2,
“Notes on Rule Fetishism and Realism”)에서의 Frank, \emph{The Bramble
Bush}에서의 Llewellyn의 논증은 다음과 같은 저작들을 배경으로 고려할 수
있다: Dickinson, “Legal Rules: Their Function in the Process of
Decision,” \emph{University of Pennsylvania Law Review} 제79권(1931);
“법 배후의 법(The Law behind the Law),” \emph{Columbia Law Review}
제29권(1929); “The Problem of the Unprovided Case,” \emph{Recueil
d'Études sur les sources de droit en l'honneur de F. Geny}, 제11권
제5장; Kantorowicz, “Some Rationalism about Realism,” \emph{Yale Law
Review} 제43권(1934).

\noteentry 139쪽. \notetitle{실망한 절대주의자로서의 회의주의자(The sceptic as a disappointed absolutist).} Miller, “Rules and Exceptions,”
\emph{International Journal of Ethics} 제66권(1956)을 보라.

\noteentry 140쪽. \notetitle{규칙의 직관적 적용(Intuitive application of rules).} Hutcheson, “The Judgement Intuitive”; “The Function of the `Hunch' in
Judicial Decision,” \emph{Cornell Law Quarterly} 제14권(1928)을 보라.

\noteentry 141쪽. \notetitle{‘헌법은 판사들이 말하는 바 그대로이다(The constitution is what the judges say it is).’} 이 표현은 미국의 Hughes 대법원장(Chief
Justice Hughes)에게서 유래한 것으로 보이며, Hendel, \emph{Charles Evan
Hughes and the Supreme Court} (1951), 제11--12쪽 참조. 그러나 Hughes
자신의 저작인 \emph{The Defence Court of the United States} (1966년판)
제37, 41쪽에서는 헌법 해석에 있어 개인의 정치적 견해로부터 독립된
판사들의 책무(duty)를 강조하고 있다.

\noteentry 149쪽. \notetitle{의회 주권에 대한 대안적 분석(Alternative analyses of the sovereignty of Parliament).} H. W. R. Wade, “The Basis of Legal
Sovereignty,” \emph{Cambridge Law Journal} (1955)을 보라. 이 견해는
Marshall, \emph{Parliamentary Sovereignty and the Commonwealth}, 제4장과
제5장에서 비판받고 있다.

\noteentry 149쪽. \notetitle{의회 주권과 신의 전능성(Parliamentary sovereignty and divine omnipotence).} Mackie, “Evil and Omnipotence,” \emph{Mind},
1955년, 제211쪽 참조.

\noteentry 150쪽. \notetitle{의회를 구속하는 것인가, 재규정하는 것인가(Binding or redefining Parliament).} 이 구분에 관해서는 Friedmann, “Trethowan's
Case, Parliamentary Sovereignty and the Limits of Legal Change,”
\emph{Australian Law Journal} 제24권(1950); Cowen, “Legislature and
Judiciary,” \emph{Modern Law Review} 제15권(1952) 및 제16권(1953);
Dixon, “The Law and the Constitution,” \emph{Law Quarterly Review}
제51권(1935); Marshall, 앞의 책, 제4장을 보라.

\noteentry 151쪽. \notetitle{1911년 및 1949년 의회법들(Parliament Acts 1911 and 1949).} 이들 법률이 위임입법(delegated legislation)의 한 형식을 수권하는 것으로
해석하는 견해는 H. W. R. Wade, 위 저작 및 Marshall, 위 저작 제44--46쪽
참조.

\noteentry 152쪽. \notetitle{1931년 웨스트민스터 제정법(Statute of Westminster), 제4조.} 권위 있는 견해들의 무게는, 이 조항의 제정이 도미니언의 동의 없이
도미니언을 위해 입법할 권한의 불가역적 종결을 구성할 수는 없었다는
견해를 지지한다. \emph{British Coal Corporation} v. \emph{The King}
(1935), AC 500; Wheare, \emph{The Statute of Westminster and Dominion
Status}, 제5판, 제297--298쪽; Marshall, 위 저작 제146--147쪽 참조. 이에
반하여, “한 번 부여된 자유는 철회될 수 없다(Freedom once conferred
cannot be revoked)”는 견해는 남아프리카 법원에서 \emph{Ndlwana} v.
\emph{Hofmeyr} (1937), AD 229, 제237쪽에서 표현된 바 있다.

\subsection{제7장 제3판
주석}\label{uxc81c7uxc7a5-uxc81c3uxd310-uxc8fcuxc11d}

\noteentry 124--128쪽. \notetitle{법의 개방적 직조(Open texture of law).} 일반적인
논의로는 Avishai Margalit, “Open Texture” (1979), \emph{Meaning and
Use} 제3권 141쪽을 보라. 하트(Hart)가 ‘개방적 직조(open texture)’ 개념을
어떻게 사용하는지에 대해서는 Brian Bix, \emph{Law, Language, and Legal
Determinacy} (Oxford University Press, 1993), 제1장을 참조하라. 법의
불확정성(indeterminacy)의 다른 원천들에 대해서는 Kelsen, \emph{Pure
Theory of Law}, 제348--356쪽을 참조. 불확정성과 그에 따른
재량(discretion)에 대한 드워킨(Dworkin)의 비판은 \emph{Taking Rights
Seriously} 제31--39쪽, 제4장, 제13장과 \emph{A Matter of Principle}
(Harvard University Press, 1986), 제5장에서 볼 수 있다. 드워킨의 이러한
비판은 David Brink, “Legal Theory, Legal Interpretation and Judicial
Review” (1988), \emph{Philosophy \& Public Affairs} 제17권 105쪽 및
Nicos Stavropoulos, “Hart's Semantics” in Jules Coleman 편,
\emph{Hart's Postscript} 중 특히 제88--98쪽에서 지지를 얻는다. 반면에
Andrei Marmor, \emph{Interpretation and Legal Theory}, 제2판 (Hart
Publishing, 2005), 제7장에서는 이 비판을 수용하지 않는다.

\noteentry 128--134쪽. \notetitle{일반 규칙들과 특정 사례들(General rules and particular cases).} 규칙 기반 결정(rule-based decision-making)의 본질에 대해서는
David Lyons, \emph{Forms and Limits of Utilitarianism} (Oxford
University Press, 1965), Frederick Schauer, \emph{Playing By the Rules:
A Philosophical Examination of Rule-Based Decision-Making in Law and in
Life} (Oxford University Press, 1991), 특히 제5장, 그리고 Larry
Alexander와 Emily Sherwin, \emph{The Rule of Rules: Morality, Rules, and
Dilemmas of Law} (Duke University Press, 2001), 제1장과 제2장을
참조하라.

\noteentry 128--129쪽. \notetitle{전범적 사례들(Paradigm cases).} 전범적 사례들이
수행하는 역할에 대한 하트의 이해를 옹호하는 논의로는 Timothy Endicott,
\emph{Vagueness in Law} (Oxford University Press, 2000), 제7장을 보라.
이에 반대하는 드워킨의 견해는 \emph{Law's Empire}, 제3장, 특히
제90--94쪽에 제시되어 있다. 이에 대한 응답으로는 Endicott, “Herbert
Hart and the Semantic Sting” in Jules Coleman 편, \emph{Hart's
Postscript} (Oxford University Press, 2001)을 보라.

\noteentry 129--130쪽. \notetitle{형식주의(Formalism).} ‘형식주의(formalism)’라는 용어는
다른 방식으로도 사용된다. 관련 논의로는 Martin Stone, “Formalism” in
Jules Coleman and Scott Shapiro 편, \emph{The Oxford Handbook of
Jurisprudence and Philosophy of Law} (Oxford University Press, 2002)을
참조하라.

\noteentry 130--132쪽. \notetitle{불확정성과 재량(Indeterminacy and discretion).} 법은,
이해 가능한 법적 질문에 대해 \emphksans{하나의 유일한(unique)} 답을
정당화하지 못할 때, 즉 법이 불완전할 때 불확정적(indeterminate)이다.
이는 어떤 법적 사례가 ‘어렵다(hard)’는 것과 동일하지 않다. 후자는 유일한
답이 존재하지만 그 답이 분명하지 않거나 이견의 대상이 되는 경우에도 참일
수 있다. 또한 이는 법이 실제로(인과적으로) 사법적 결정들을 결정짓는다는
견해와도 다르다. 법은 전적으로 확정적일 수 있지만 판사들이 그것을 알지
못하거나, 알고도 적용하지 않기로 결정할 수 있다. 법의 확정성(legal
determinacy)은 결정의 예측가능성(decisional predictability)을 함의하지
않는다. (비교: Brian Leiter, “Legal Indeterminacy” (1995), \emph{Legal
Theory} 제1권 481쪽) 하트가 옹호하는 것은 법의 정당화적
불확정성(justificatory indeterminacy)이다. 하트의 테제에 대한 거부는
Ronald Dworkin, “No Right Answer?” in P. M. S. Hacker and Joseph Raz
편, \emph{Law, Morality and Society}, 그리고 Ronald Dworkin, “On Gaps
in the Law” in Paul Amselek and Neil MacCormick 편, \emph{Controversies
About Law's Ontology} (Edinburgh University Press, 1991)에서 볼 수 있다.
이에 대한 옹호는 Joseph Raz, \emph{The Authority of Law}, 제4장에서 찾을
수 있다.

\noteentry 131--133쪽. \notetitle{규칙들과 가변적 기준들(Rules and variable standards).} 이 지점에서 하트(Hart)는 규칙(rule)을 두어야 하는지
\emphksans{여부(whether)}를 논의하는 것이 아니라, 다양한 상황에 어떤
\emphksans{종류(kind)}의 규칙이 적절할지를 논의하고 있다. 미국의 법률가들은,
헨리 하트(Henry Hart, H. L. A. Hart와는 무관)와 앨버트 삭스(Albert
Sacks)의 영향을 받아, 때때로 ‘규칙들(rules)’과 ‘기준들(standards)’을
구분하지만, H. L. A. Hart에게 기준들(standards)은 규칙들의 한 종류이다.
이 구분에 대해서는 Henry M. Hart and Albert Sacks, \emph{The Legal
Process: Basic Problems in the Making and Application of Law} (W. N.
Eskridge, Jr.~and P. P. Frickey 편집, Foundation Press, 1994)
139--141쪽, 그리고 Pierre J. Schlag, “Rules and Standards” (1985)
\emph{UCLA Law Review} 제33권 379쪽을 참조하라. 마찬가지로, 흔히
‘원칙들(principles)’이라 불리는 것들도 하트의 규칙들에 관한 관념(idea of
rules)에 포함되며, 그는 드워킨(Dworkin)이 \emph{Taking Rights Seriously}
제22--28쪽에서 설정한 것과 같은 규칙과 원칙 간의 범주적 구분을 거부한다.
이에 대한 하트의 입장은 \emph{Postscript} 제261--263쪽에 언급되어 있다.
관련 논의로는 Joseph Raz, “Legal Principles and the Limits of Law”
(1972) \emph{Yale Law Journal} 제81권 823쪽, Larry Alexander and Ken
Kress, “Against Legal Principles” (1997) \emph{Iowa Law Review} 제82권
739쪽을 참조하라.

\noteentry 134--135쪽. \notetitle{선례(Precedent).} 이 주제에 관한 문헌은 매우 풍부하다.
참고문헌으로는 다음을 보라: A. W. B. Simpson, “The Common Law and Legal
Theory” in A. W. B. Simpson 편, \emph{Oxford Essays in Jurisprudence};
Ronald Dworkin, \emph{Taking Rights Seriously} 제110--123쪽; Neil
MacCormick, \emph{Legal Reasoning and Legal Theory} (수정판, Oxford
University Press, 1994); Laurence Goldstein 편, \emph{Precedent in Law}
(Oxford University Press, 1987); Stephen Perry, “Judicial Obligation,
Precedent and the Common Law” (1987) \emph{Oxford Journal of Legal
Studies} 제7권 215쪽; Frederick Schauer, “Precedent” (1987)
\emph{Stanford Law Review} 제39권 571쪽; Susan Hurley, “Coherence,
Hypothetical Cases, and Precedent” (1990) \emph{Oxford Journal of Legal
Studies} 제10권 221쪽; Larry Alexander, “Precedent” in Dennis
Patterson 편, \emph{A Companion to the Philosophy of Law and Legal
Theory} (Blackwell, 1996); Neil Duxbury, \emph{The Nature and Authority
of Precedent} (Cambridge University Press, 2008).

\noteentry 136--141쪽. \notetitle{규칙회의주의의 다양상들(Varieties of rule scepticism).} 하트의 관련 에세이로는 “American Jurisprudence through
English Eyes: The Nightmare and the Noble Dream” (\emph{Essays in
Jurisprudence and Philosophy}, 제4장)이 있다. 하트의 저작 이후 등장한
가장 저명한 규칙회의주의자(rule-sceptics)는 미국의 ‘비판법학(Critical
Legal Studies, CLS)’ 운동의 법학자들이었다. 철학적으로는 덜 정교하지만
영향력 있는 논의로는 Roberto Unger, \emph{The Critical Legal Studies
Movement} (Harvard University Press, 1986), 특히 제1--14쪽, 그리고
Duncan Kennedy, \emph{A Critique of Adjudication {[}fin de siècle{]}}
(Harvard University Press, 1997)를 들 수 있다. 초기 CLS 문헌에 대한
비판적 논의는 Ken Kress, “Legal Indeterminacy” (1989) \emph{California
Law Review} 제77권 283쪽을 참조하라. 또한 John Finnis, “On the Critical
Legal Studies Movement”, 그의 \emph{Philosophy of Law} 제13장, Andrew
Altman, \emph{Critical Legal Studies: A Liberal Critique} (Princeton
University Press, 1990)도 보라. Brian Leiter는 “Legal Realism and Legal
Positivism Reconsidered”, 그의 \emph{Naturalizing Jurisprudence}
(Oxford University Press, 2007) 제2장에서 하트의 규칙회의주의 다양상들을
다룬다.
}
